% ROZDZIAL 6 -- Semantyka przenoszenia, kategorie wartosci, perfect forwarding (C++20)
% Fragment do wstawienia za rozdzialem 5.
% Wymagana preambula: pakiet listings i konfiguracja lstlisting.
% Uzycie: % ROZDZIAL 6 -- Semantyka przenoszenia, kategorie wartosci, perfect forwarding (C++20)
% Fragment do wstawienia za rozdzialem 5.
% Wymagana preambula: pakiet listings i konfiguracja lstlisting.
% Uzycie: % ROZDZIAL 6 -- Semantyka przenoszenia, kategorie wartosci, perfect forwarding (C++20)
% Fragment do wstawienia za rozdzialem 5.
% Wymagana preambula: pakiet listings i konfiguracja lstlisting.
% Uzycie: % ROZDZIAL 6 -- Semantyka przenoszenia, kategorie wartosci, perfect forwarding (C++20)
% Fragment do wstawienia za rozdzialem 5.
% Wymagana preambula: pakiet listings i konfiguracja lstlisting.
% Uzycie: \input{rozdzial_6_move_semantics_cpp.tex}

\chapter[Semantyka przenoszenia w C++]{Semantyka przenoszenia, kategorie wartości i przekazywanie doskonałe}
\label{chap:move-semantics-cpp}

\section{Dlaczego wprowadzono semantykę przenoszenia?}

Rozważmy kontener przechowujący milion liczb. Skopiowanie go oznacza
zwykle przydzielenie nowej pamięci i skopiowanie każdego elementu.
Jeżeli dotychczasowy obiekt nie będzie już potrzebny, bardziej naturalna
jest operacja \emph{przekazania własności} jego zasobu zamiast kopiowania.
Od C++11 język udostępnia \emph{semantykę przenoszenia} (\textit{move
semantics}): programista może rozróżnić sytuację, w której obiekt źródłowy
ma pozostać niezmieniony, od sytuacji, w której jego zasoby wolno wykorzystać
ponownie.

\begin{lstlisting}
#include <string>
#include <utility>
#include <vector>

std::vector<std::string> a(100'000, "data");
auto b = a;                // copy: a remains intact

auto c = std::move(a);     // move may transfer a's allocation
// c now owns its elements; a is still valid, but its state
// after this move should not be assumed without a guarantee.
\end{lstlisting}

Nie oznacza to, że każde przeniesienie jest szybsze niż kopiowanie.
Dla typu \texttt{int} oba działania są w praktyce tożsame, a niektóre typy
kopiują część małych danych również podczas przenoszenia.
Kluczowa idea jest semantyczna: \emph{przeniesienie pozwala zmienić stan
źródła, a kopiowanie nie powinno go zmieniać}.

\paragraph{Zakres rozdziału.}
Zakładamy \textbf{C++20}. Kategorie wartości istnieją od początku języka,
referencje do r-wartości i \texttt{std::move} wprowadzono w C++11,
\emph{guaranteed copy elision} w C++17, a wiele przykładów poniżej
wykorzystuje koncepcje omówione w rozdziale~5.
Wszystkie przykłady zawierają komentarze angielskie, aby łatwo
przenosić je do projektów.

\section{Typ wyrażenia a kategoria wartości}
\label{sec:move-categories}

\textbf{Typ} i \textbf{kategoria wartości} są dwiema różnymi własnościami
wyrażenia. Typ mówi, jaki rodzaj danych otrzymujemy; kategoria mówi
m.in., czy wyrażenie identyfikuje istniejący obiekt i jak można
wiązać je z referencjami. Dokładnie trzy podstawowe kategorie to:
\emph{lvalue}, \emph{xvalue} i \emph{prvalue}.

\begin{center}
\begin{tabular}{|l|p{10.5cm}|}
\hline
\textbf{Kategoria} & \textbf{Intuicja i przykłady} \\
\hline
\texttt{lvalue} & Wyrażenie identyfikujące obiekt lub funkcję; np. \texttt{x}, \texttt{*ptr}, \texttt{++x} dla wbudowanego \texttt{int}. \\
\hline
\texttt{xvalue} & \emph{Expiring value}: wyrażenie identyfikujące obiekt, którego zasoby mogą być ponownie wykorzystane; np. \texttt{std::move(x)}. \\
\hline
\texttt{prvalue} & \emph{Pure rvalue}: wyrażenie obliczające wartość lub inicjalizujące obiekt; np. \texttt{42}, \texttt{x + 1}, \texttt{std::string("hi")}. \\
\hline
\end{tabular}
\end{center}

Pozostałe dwie nazwy oznaczają sumy kategorii:
\[
\underbrace{\mathrm{glvalue}}_{\mathrm{lvalue}\,\cup\,\mathrm{xvalue}},
\qquad
\underbrace{\mathrm{rvalue}}_{\mathrm{prvalue}\,\cup\,\mathrm{xvalue}}.
\]
Nie należy utożsamiać \texttt{rvalue} z tymczasowym obiektem ani
\texttt{lvalue} z wyrażeniem znajdującym się po lewej stronie przypisania.
Są to techniczne kategorie zdefiniowane przez język.

\subsection{Kilka wyrażeń o zaskakujących kategoriach}

\begin{lstlisting}
#include <string>
#include <utility>

int x = 7;
int* p = &x;

// x                 : lvalue
// *p                : lvalue
// (x)               : lvalue
// 7                 : prvalue
// x + 1             : prvalue
// std::move(x)      : xvalue
// std::string{"abc"}: prvalue

int& lref = x;
int&& rref = 9;

// lref              : lvalue
// rref              : lvalue (!) - it has a name
// std::move(rref)   : xvalue
\end{lstlisting}

To fundamentalna zasada: \textbf{nazwa zmiennej jest wyrażeniem lvalue,
nawet jeśli typ tej zmiennej to referencja rvalue \texttt{T\&\&}}.
Samo zadeklarowanie \texttt{T\&\&} nie oznacza, że każde użycie takiej
zmiennej będzie r-wartością.

\subsection{Badanie kategorii za pomocą \texttt{decltype((expr))}}

Dla prostego, nieujętego w nawiasy identyfikatora składnia
\texttt{decltype(x)} zwraca deklarowany typ. Dla ogólnego wyrażenia
\texttt{decltype((expr))} oddaje jego kategorię poprzez typ:
\[
\mathrm{decltype}((e))=
\left\{
\begin{array}{ll}
T\&, & \mbox{gdy } e \mbox{ jest lvalue typu } T,\\
T\&\&, & \mbox{gdy } e \mbox{ jest xvalue typu } T,\\
T, & \mbox{gdy } e \mbox{ jest prvalue typu } T.
\end{array}
\right.
\]

\begin{lstlisting}
#include <type_traits>
#include <utility>

int x = 10;
static_assert(std::is_same_v<decltype(x), int>);
static_assert(std::is_same_v<decltype((x)), int&>);
static_assert(std::is_same_v<decltype((std::move(x))), int&&>);
static_assert(std::is_same_v<decltype((x + 1)), int>);
\end{lstlisting}

\paragraph{Uwaga.} \texttt{decltype} stosuje szczególne reguły dla
nieujętych w nawiasy nazw i dostępu do składowych. Dlatego
\texttt{decltype(x)} oraz \texttt{decltype((x))} mogą zwrócić
różne typy. Nie jest to sprzeczność, lecz zamierzona cecha języka.

\subsection{Różnica między \texttt{auto} i \texttt{decltype(auto)}}

Zwykłe \texttt{auto} w deklaracji zmiennej najczęściej dedukuje typ
przez wartość i pomija referencję wynikającą z inicjalizatora.
Natomiast \texttt{decltype(auto)} stosuje reguły \texttt{decltype}
i może zachować kategorię wartości wyrażenia.

\begin{lstlisting}
#include <type_traits>

int n = 5;
auto copy = (n);           // int: independent value
decltype(auto) alias = (n); // int&: refers to n

static_assert(std::is_same_v<decltype(copy), int>);
static_assert(std::is_same_v<decltype(alias), int&>);

auto&& first = n;          // deduces int&
auto&& second = 5;         // deduces int&&
static_assert(std::is_same_v<decltype(first), int&>);
static_assert(std::is_same_v<decltype(second), int&&>);
\end{lstlisting}

W typie wyniku funkcji zapis \texttt{decltype(auto)} wymaga
ostrożności. Instrukcja \texttt{return (x);} może zwrócić
referencję do zmiennej lokalnej, która przestanie istnieć.
Zwykłe \texttt{auto} jest często bezpieczniejszym wyborem,
gdy funkcja ma zwrócić własny obiekt przez wartość.

\section{Referencje lvalue i rvalue}
\label{sec:move-references}

\begin{lstlisting}
#include <string>

std::string name = "Ada";
std::string& a = name;            // binds to mutable lvalue
const std::string& b = name;      // binds to lvalue
const std::string& c = "temporary"; // can bind to temporary
std::string&& d = std::string{"x"}; // binds to rvalue

// std::string& e = std::string{"x"}; // error
// std::string&& f = name;           // error
\end{lstlisting}

Referencja \texttt{T\&} umożliwia pracę z istniejącym, zwykle
nazwanym obiektem; \texttt{T\&\&} może wiązać się z r-wartością.
\texttt{const T\&} wiąże się zarówno z lvalue, jak i z rvalue,
nie dając prawa do modyfikowania obiektu poprzez tę referencję.

\subsection{Wybór przeciążenia}

\begin{lstlisting}
#include <iostream>
#include <string>
#include <utility>

void inspect(const std::string&) {
    std::cout << "const lvalue reference\n";
}
void inspect(std::string&&) {
    std::cout << "rvalue reference\n";
}

int main() {
    std::string text = "hello";
    inspect(text);             // const lvalue reference
    inspect(std::string{"x"}); // rvalue reference
    inspect(std::move(text));  // rvalue reference
}
\end{lstlisting}

Istnienie przeciążenia \texttt{f(T\&\&)} samo w sobie nie przesądza
o rzeczywistym przeniesieniu czegokolwiek. Funkcja może tylko
odczytać argument; \emph{move} jest operacją wykonywaną przez
konstruktor lub przypisanie, a nie przez mechanizm wiązania referencji.

\subsection{Wydłużenie życia obiektu tymczasowego}

W typowych deklaracjach lokalnych związanie obiektu tymczasowego
bezpośrednio z referencją może wydłużyć jego życie do końca życia
tej referencji:
\begin{lstlisting}
#include <string>

const std::string& a = std::string{"safe until a ends"};
std::string&& b = std::string{"safe until b ends"};
\end{lstlisting}
Nie wolno jednak uogólniać tej reguły na zwracanie referencji ani
na dowolne przekazywanie referencji dalej. Przykładowo:
\begin{lstlisting}
#include <string>

const std::string& broken() {
    return std::string{"dangling"}; // dangling reference
}
\end{lstlisting}
Zwrócona referencja nie utrzymuje tego obiektu przy życiu.
Jeżeli funkcja ma \emph{wytworzyć} wynik, zwykle powinna zwracać go
przez wartość.

\section{Co naprawdę robi \texttt{std::move}?}
\label{sec:move-std-move}

Nazwa jest nieco myląca: \texttt{std::move(x)} \textbf{nie przenosi
zasobów samodzielnie}. W przybliżeniu wykonuje rzutowanie do
referencji rvalue odpowiedniego typu. Dopiero dalszy konstruktor,
operator przypisania lub inna funkcja może potraktować wynik
jako obiekt nadający się do przeniesienia.

\begin{lstlisting}
#include <type_traits>
#include <utility>

template<class T>
constexpr std::remove_reference_t<T>&&
my_move(T&& value) noexcept {
    return static_cast<std::remove_reference_t<T>&&>(value);
}
\end{lstlisting}

Powyższa implementacja ilustruje istotę działania
\texttt{std::move}. W rzeczywistym kodzie korzystamy z funkcji
standardowej z nagłówka \texttt{<utility>}.

\begin{lstlisting}
#include <string>
#include <utility>

std::string first = "long text";
std::string second = std::move(first);
// A move constructor is selected for second.
// first is valid; do not assume its exact contents.

first = "reused";             // valid: assignment
first.clear();                // valid: clear()
\end{lstlisting}

\paragraph{Stan obiektu po przeniesieniu.}
Typy biblioteczne zwykle gwarantują stan \emph{poprawny, lecz
nieokreślony co do zawartości} (\textit{valid but unspecified}), o ile
kontrakt konkretnego typu nie obiecuje więcej. Takiego obiektu
można bezpiecznie zniszczyć, przypisać mu nową wartość albo
wykonać operację, której wymagania wstępne na to pozwalają.
Nie należy jednak \emph{zakładać}, że każdy przeniesiony
\texttt{std::string} jest pusty. Dla niektórych typów, np.
\texttt{std::unique\_ptr}, stan po przeniesieniu określono dokładniej.

\subsection{Problem obiektu \texttt{const}}

\begin{lstlisting}
#include <string>
#include <utility>

const std::string original = "cannot be emptied here";
std::string other = std::move(original);
// Typically: copy, not move (const string&& cannot bind to
// string's non-const move constructor).
\end{lstlisting}

\texttt{std::move(original)} ma w tym przypadku typ
\texttt{const std::string\&\&}. Typowy konstruktor przenoszący
przyjmuje \texttt{std::string\&\&}, więc nie może otrzymać
argumentu \texttt{const}. Kompilator wybiera konstruktor kopiujący
przyjmujący \texttt{const std::string\&}. \textbf{Samo wystąpienie
\texttt{std::move} w kodzie nie gwarantuje przeniesienia}.\footnote{Dla
typów niekopiowalnych ten sam zapis może skończyć się błędem
kompilacji.}

\section{Konstruktor przenoszący i przenoszący operator przypisania}
\label{sec:move-special-members}

W klasie \texttt{T} konstruktor przenoszący ma zwykle sygnaturę
\texttt{T(T\&\&)}, natomiast przenoszący operator przypisania:
\texttt{T\& operator=(T\&\&)}. Nie należy mylić ich z kopiowaniem:
\begin{lstlisting}
struct Item {
    Item() = default;
    Item(const Item&);            // copy constructor
    Item& operator=(const Item&); // copy assignment
    Item(Item&&) noexcept;        // move constructor
    Item& operator=(Item&&) noexcept; // move assignment
    ~Item();                      // destructor
};
\end{lstlisting}

\subsection{Przykład z obserwacją operacji}
\label{sec:move-trace}

\begin{lstlisting}
#include <iostream>
#include <utility>

struct Trace {
    Trace() { std::cout << "default\n"; }
    Trace(const Trace&) { std::cout << "copy ctor\n"; }
    Trace(Trace&&) noexcept { std::cout << "move ctor\n"; }
    Trace& operator=(const Trace&) {
        std::cout << "copy assign\n";
        return *this;
    }
    Trace& operator=(Trace&&) noexcept {
        std::cout << "move assign\n";
        return *this;
    }
};

int main() {
    Trace a;                    // default
    Trace b = a;                // copy ctor
    Trace c = std::move(a);     // move ctor
    b = c;                      // copy assign
    b = std::move(c);           // move assign
}
\end{lstlisting}

Konstruktor tworzy \emph{nowy} obiekt, a operator przypisania
zmienia już istniejący. W tym przykładzie \texttt{Trace} jedynie
wypisuje komunikaty. W klasie zarządzającej zasobem konstruktor
przenoszący powinien przekazać własność zasobu, a operator
przypisania musi dodatkowo poprawnie pozbyć się zasobu
posiadanego wcześniej przez cel.

\subsection{Pełny przykład: klasa samodzielnie zarządzająca tablicą}
\label{sec:move-raw-buffer}

To przykład \emph{edukacyjny}. W rzeczywistym projekcie zwykle
lepiej użyć standardowego kontenera, np. \texttt{std::vector<int>}.

\begin{lstlisting}
#include <algorithm>
#include <cstddef>
#include <utility>

class RawBuffer {
    std::size_t size_ = 0;
    int* data_ = nullptr;

public:
    explicit RawBuffer(std::size_t n)
        : size_(n), data_(n ? new int[n]{} : nullptr) {}

    ~RawBuffer() { delete[] data_; }

    RawBuffer(const RawBuffer& other)
        : RawBuffer(other.size_) {
        if (size_ != 0)
            std::copy_n(other.data_, size_, data_);
    }

    RawBuffer& operator=(const RawBuffer& other) {
        if (this != &other) {
            RawBuffer copy(other); // may throw; *this unchanged
            swap(copy);
        }
        return *this;
    }

    RawBuffer(RawBuffer&& other) noexcept
        : size_(std::exchange(other.size_, 0)),
          data_(std::exchange(other.data_, nullptr)) {}

    RawBuffer& operator=(RawBuffer&& other) noexcept {
        if (this != &other) {
            delete[] data_;
            size_ = std::exchange(other.size_, 0);
            data_ = std::exchange(other.data_, nullptr);
        }
        return *this;
    }

    void swap(RawBuffer& other) noexcept {
        using std::swap;
        swap(size_, other.size_);
        swap(data_, other.data_);
    }

    std::size_t size() const noexcept { return size_; }
    int& operator[](std::size_t i) { return data_[i]; }
    const int& operator[](std::size_t i) const { return data_[i]; }
};
\end{lstlisting}

\texttt{std::exchange(obj, value)} zwraca starą wartość \texttt{obj}
i zastępuje ją wartością \texttt{value}. Konstruktor przenoszący
przejmuje więc wskaźnik i rozmiar, a źródło ustawia w jawnie
określony stan pusty. Operator przenoszącego przypisania
usuwa dotychczasową tablicę celu, przejmuje nową i zabezpiecza
przypadek \texttt{b = std::move(b)}. W implementacji kopiowania
widać z kolei \emph{silną gwarancję wyjątkową}: jeśli utworzenie
kopii nie powiedzie się, pierwotny obiekt pozostaje bez zmian.

\section{Reguła pięciu, reguła trzech i reguła zera}
\label{sec:move-rule-of-five}

Klasy zarządzające zasobem często muszą kontrolować pięć operacji:
(1) destruktor, (2) konstruktor kopiujący, (3) kopiujący operator
przypisania, (4) konstruktor przenoszący oraz (5) przenoszący
operator przypisania. To \emph{reguła pięciu}.

Historyczna \emph{reguła trzech} odnosiła się do destruktora,
konstruktora kopiującego i kopiującego operatora przypisania,
zanim język udostępnił move semantics.

\paragraph{Ważne: deklarowanie destruktora może blokować generowanie
konstruktora przenoszącego.}

\begin{lstlisting}
#include <memory>
#include <type_traits>

struct WithDestructor {
    std::unique_ptr<int> p;
    ~WithDestructor() = default;
};

static_assert(!std::is_move_constructible_v<WithDestructor>);

struct WithMove {
    std::unique_ptr<int> p;
    ~WithMove() = default;
    WithMove() = default;
    WithMove(WithMove&&) noexcept = default;
    WithMove& operator=(WithMove&&) noexcept = default;
};
static_assert(std::is_move_constructible_v<WithMove>);
\end{lstlisting}

W \texttt{WithDestructor} użytkownik zadeklarował destruktor,
więc kompilator nie deklaruje automatycznie operacji przenoszących.
Kopiowanie jest niedostępne z powodu \texttt{unique\_ptr}.
Samo \texttt{= default} przy destruktorze nie usuwa tej reguły.

\subsection{Najlepsza praktyka: reguła zera}

\emph{Reguła zera} mówi, że klasy, które nie są bezpośrednio
odpowiedzialne za posiadanie surowych zasobów, zwykle nie powinny
same definiować żadnej z pięciu operacji.

\begin{lstlisting}
#include <string>
#include <vector>

struct Dataset {
    std::string title;
    std::vector<double> values;
};

Dataset original{"results", {1.0, 2.0, 3.0}};
Dataset copy = original;           // deep copy of vector data
Dataset moved = std::move(original); // member-wise move
\end{lstlisting}

To kontenery standardowe odpowiadają za przydział i zwalnianie
pamięci. Automatycznie generowane operacje \texttt{Dataset}
wykonują odpowiednie operacje na jego składowych. \textbf{Zalecenie:}
zacznij od reguły zera; własny konstruktor przenoszący pisz
jedynie wtedy, gdy zachowanie domyślne nie jest właściwe.

\section{Przenoszenie a wyjątki: \texttt{noexcept}}
\label{sec:move-noexcept}

Kontenery takie jak \texttt{std::vector<T>} przy realokacji muszą
przenieść lub skopiować dotychczasowe elementy do nowego obszaru.
Jeśli przenoszenie \texttt{T} może rzucać wyjątki, a \texttt{T}
można kopiować, biblioteka może wybrać kopiowanie, by zachować
silniejsze gwarancje wyjątkowe. Właśnie dlatego poprawne operacje
przenoszące, które nie mogą rzucać wyjątków, warto deklarować
\texttt{noexcept}.

\begin{lstlisting}
#include <type_traits>
#include <utility>

struct Resource {
    Resource() = default;
    Resource(const Resource&) = default;
    Resource(Resource&&) noexcept = default;
};

static_assert(std::is_nothrow_move_constructible_v<Resource>);

// move_if_noexcept can provide an lvalue reference if
// moving might throw and a copy is available.
\end{lstlisting}

\texttt{std::move\_if\_noexcept(x)} zwraca odpowiednio
\texttt{T\&\&} albo \texttt{const T\&}, zależnie od własności
konstruktorów danego typu. To narzędzie używane przy implementowaniu
kontenerów, a nie uniwersalny zamiennik \texttt{std::move}.

\paragraph{Ostrożnie z \texttt{noexcept}.}
Wyjście wyjątku poza funkcję zadeklarowaną jako
\texttt{noexcept} wywołuje \texttt{std::terminate}.
Dlatego bezwyjątkowość musi wynikać z rzeczywistych operacji składowych.

\section{Zwijanie referencji (\textit{reference collapsing})}
\label{sec:move-ref-collapse}

Gdy referencje występują pośrednio poprzez aliasy lub
podstawienie argumentów szablonu, stosuje się \emph{zwijanie
referencji}. Wszystkie przypadki przedstawia tabela:

\begin{center}
\begin{tabular}{|c|c|c|}
\hline
\textbf{Typ pośredni} & \textbf{Zapis} & \textbf{Wynik} \\
\hline
\texttt{T\&} i \texttt{\&} & \texttt{T\& \&} & \texttt{T\&} \\
\texttt{T\&} i \texttt{\&\&} & \texttt{T\& \&\&} & \texttt{T\&} \\
\texttt{T\&\&} i \texttt{\&} & \texttt{T\&\& \&} & \texttt{T\&} \\
\texttt{T\&\&} i \texttt{\&\&} & \texttt{T\&\& \&\&} & \texttt{T\&\&} \\
\hline
\end{tabular}
\end{center}

Krótko: jeżeli \emph{choć jeden} składnik jest \texttt{\&},
wynik jest \texttt{\&}; tylko \texttt{\&\&} połączone
z \texttt{\&\&} daje \texttt{\&\&}.

\begin{lstlisting}
#include <type_traits>

using L = int&;
using R = int&&;

static_assert(std::is_same_v<L&, int&>);
static_assert(std::is_same_v<L&&, int&>);
static_assert(std::is_same_v<R&, int&>);
static_assert(std::is_same_v<R&&, int&&>);
\end{lstlisting}

Zwijanie referencji jest podstawą \emph{perfect forwarding}:
dzięki niemu jeden szablon może przyjmować zarówno lvalue,
jak i rvalue i przekazywać tę różnicę dalej.

\section{Referencje przekazujące i dedukcja \texttt{T\&\&}}
\label{sec:move-forwarding-ref}

Zapis \texttt{T\&\&} ma dwa znaczenia w zależności od kontekstu.
Jeśli \texttt{T} jest dedukowanym, niekwalifikowanym parametrem
typu \emph{tego samego} szablonu funkcji, \texttt{T\&\&}
to \emph{referencja przekazująca} (\textit{forwarding reference}).
W pozostałych typowych kontekstach jest to po prostu
referencja rvalue.

\begin{lstlisting}
#include <string>

template<class T>
void forwarding(T&& x);         // forwarding reference

template<class T>
void not_forwarding(const T&& x); // rvalue reference

void fixed(std::string&& x);    // rvalue reference

template<class T>
struct Holder {
    void receive(T&& x);        // NOT forwarding: T comes
                               // from the class template

    template<class U>
    void relay(U&& x);          // forwarding: U is deduced
};
\end{lstlisting}

Dla wywołania \texttt{forwarding(x)} z lvalue typu \texttt{int}
kompilator dedukuje \texttt{T = int\&}, zatem parametr ma po
zwinięciu typ \texttt{int\&}. Dla \texttt{forwarding(5)}
dedukuje \texttt{T = int}, a parametr ma typ \texttt{int\&\&}.

\begin{lstlisting}
#include <type_traits>

template<class T>
void check(T&& x) {
    if constexpr (std::is_lvalue_reference_v<T>) {
        // Called with an lvalue.
    } else {
        // Called with an rvalue.
    }
}
\end{lstlisting}

\paragraph{Drobne wyjątki.} \texttt{auto\&\&} zwykle również
pozwala na przekazywanie kategorii, lecz dedukcja \texttt{auto}
ma szczególne reguły dla list inicjalizacyjnych.
\texttt{const T\&\&} nie jest referencją przekazującą.

\section{Perfect forwarding i \texttt{std::forward}}
\label{sec:move-perfect-forwarding}

Przez \emph{perfect forwarding} rozumiemy przekazanie argumentu
do następnej funkcji przy zachowaniu kategorii wartości,
a także odpowiednich kwalifikatorów typu.

\begin{lstlisting}
#include <iostream>
#include <utility>

void target(int&)  { std::cout << "lvalue\n"; }
void target(int&&) { std::cout << "rvalue\n"; }

template<class T>
void incorrect(T&& value) {
    target(value); // value is a named variable: always lvalue
}

template<class T>
void correct(T&& value) {
    target(std::forward<T>(value));
}

int main() {
    int x = 1;
    incorrect(2); // lvalue (!)
    correct(x);   // lvalue
    correct(2);   // rvalue
}
\end{lstlisting}

W przybliżeniu \texttt{std::forward<T>(value)} zwraca
lvalue, jeżeli \texttt{T} wydedukowano jako referencję lvalue,
i xvalue w przeciwnym razie. \texttt{std::move(value)} zawsze
wymusza kategorię xvalue: w funkcji przekazującej byłoby to
błędem, ponieważ niepotrzebnie pozwalałoby na przeniesienie
zasobu z argumentu przekazanego jako lvalue.

\subsection{Przykład: fabryka obiektów}

\begin{lstlisting}
#include <memory>
#include <string>
#include <utility>

struct Person {
    std::string name;
    int age;

    Person(std::string n, int a)
        : name(std::move(n)), age(a) {}
};

template<class T, class... Args>
std::unique_ptr<T> make_object(Args&&... args) {
    return std::unique_ptr<T>(
        new T(std::forward<Args>(args)...)
    );
}

int main() {
    std::string first = "Ada";
    auto p = make_object<Person>(first, 37); // first copied
    auto q = make_object<Person>(std::move(first), 38);
}
\end{lstlisting}

To edukacyjna wersja fabryki. W produkcyjnym kodzie
\texttt{std::make\_unique<T>(...)} robi to samo w prostszy
sposób. W drugim wywołaniu użytkownik świadomie pozwolił na
przeniesienie \texttt{first}; w pierwszym pozostaje on
nienaruszony. \texttt{std::forward<Args>(args)...} rozwija pakiet
argumentów omówiony w rozdziale~4.

\subsection{Nie przekazuj tego samego zasobu dwa razy bez zastanowienia}

\begin{lstlisting}
#include <utility>

template<class T>
void bad_double_forward(T&& x) {
    consume(std::forward<T>(x));
    consume(std::forward<T>(x)); // may use moved-from object
}
\end{lstlisting}

Ten fragment jest \emph{schematem błędu}: zakładamy, że istnieje
funkcja \texttt{consume}. Jeśli pierwsze wywołanie wykorzysta
argument rvalue przez przeniesienie, drugie otrzyma już obiekt
po przeniesieniu. Referencję przekazującą traktuj zatem jako
przeznaczoną zwykle do \emph{jednego} przekazania dalej.

\section{Perfect forwarding a przeciążenia i szablony klas}

Konstruktor z referencją przekazującą potrafi być bardzo zachłanny
podczas rozstrzygania przeciążeń. Może np. wygrać z konstruktorem
kopiującym przy argumentach nie-\texttt{const}.

\begin{lstlisting}
#include <string>
#include <utility>

struct Named {
    std::string name;

    Named(const Named&) = default;

    template<class U>
    explicit Named(U&& u)
        : name(std::forward<U>(u)) {}
};

// Named a{"text"};
// Named b(a);  // may select the constructor template,
//              // then fail to initialize std::string from Named&
\end{lstlisting}

Zwykle lepiej zacząć od zwykłych przeciążeń albo konstruktora
przyjmującego konkretny typ. Jeżeli szablon jest potrzebny,
można ograniczyć go koncepcją:

\begin{lstlisting}
#include <concepts>
#include <string>
#include <type_traits>
#include <utility>

struct SafeNamed {
    std::string name;

    SafeNamed(const SafeNamed&) = default;
    SafeNamed(SafeNamed&&) = default;

    template<class U>
      requires (!std::same_as<std::remove_cvref_t<U>, SafeNamed>
                && std::constructible_from<std::string, U&&>)
    explicit SafeNamed(U&& u)
        : name(std::forward<U>(u)) {}
};
\end{lstlisting}

Takie powiązanie rozdziału o perfect forwarding z koncepcjami
z rozdziału~5 jest bardzo przydatne
w bibliotekach generycznych.

\section{Copy elision: kiedy nie ma ani kopii, ani przeniesienia?}
\label{sec:move-elision}

Prosty model \emph{zwrócenie przez wartość = kopia lub move}
jest niepoprawny w nowoczesnym C++. Od C++17 pewne wyrażenia
\texttt{prvalue} inicjalizują wynik \emph{bezpośrednio w miejscu
docelowym}. Nie istnieje dodatkowy obiekt, który trzeba by
skopiować lub przenieść.

\begin{lstlisting}
struct NonMovable {
    NonMovable() = default;
    NonMovable(const NonMovable&) = delete;
    NonMovable(NonMovable&&) = delete;
};

NonMovable create() {
    return NonMovable{}; // guaranteed elision, C++17+
}

int main() {
    NonMovable x = create(); // OK in C++17+
}
\end{lstlisting}

Nie jest to zwykła optymalizacja, którą kompilator \emph{może}
wyłączyć: w pokazanym przypadku wynika to z semantyki języka.

\subsection{NRVO a gwarantowana elizja}

\begin{lstlisting}
#include <string>

std::string first() {
    return std::string{"hello"}; // mandatory direct construction
}

std::string second() {
    std::string s = "hello";
    return s;                    // NRVO may occur
}
\end{lstlisting}

\textit{Named Return Value Optimization} (NRVO) dotyczy
zwracania nazwanej lokalnej zmiennej. W C++20 kompilator \emph{może}
pominąć utworzenie dodatkowego obiektu. Jeśli NRVO nie nastąpi,
wyrażenie \texttt{s} w instrukcji \texttt{return} może
uczestniczyć w niejawnej próbie przeniesienia.

\paragraph{Nie pisz automatycznie \texttt{return std::move(s);}.}

\begin{lstlisting}
std::string make_text() {
    std::string s = "result";
    return s;             // allows NRVO
}

std::string less_good() {
    std::string s = "result";
    return std::move(s);  // inhibits NRVO; move may be required
}
\end{lstlisting}

\texttt{std::move(s)} zmienia wynik w xvalue i tym samym
nie spełnia warunków NRVO dotyczących nazwy obiektu.
Jeśli funkcja zwraca własną lokalną zmienną przez wartość,
zazwyczaj najlepsze jest \texttt{return s;}.

\section{Argumenty funkcji: kiedy kopiować, kiedy przenosić?}
\label{sec:move-api}

Wybór sposobu przyjmowania argumentu jest częścią kontraktu API.
Nie chodzi jedynie o wydajność, ale przede wszystkim o własność
i intencję.

\begin{center}
\begin{tabular}{|p{3.4cm}|p{4.3cm}|p{4.7cm}|}
\hline
\textbf{Parametr} & \textbf{Znaczenie} & \textbf{Kiedy używać} \\
\hline
\texttt{const T\&} & Odczyt bez przejmowania własności. & Duże obiekty, gdy wystarczy tylko odczyt. \\
\hline
\texttt{T\&} & Modyfikacja obiektu należącego do wywołującego. & Parametry wyjściowe i operacje in-place. \\
\hline
\texttt{T} & Własna wartość; może być skopiowana albo przeniesiona. & Małe typy i funkcje mające zachować wartość. \\
\hline
\texttt{T\&\&} & Dostęp do argumentu będącego r-wartością. & Przeciążenia zależne od kategorii oraz operacje konsumujące. \\
\hline
\texttt{T\&\&} z dedukcją \texttt{T} & Referencja przekazująca. & Opakowania i fabryki zachowujące kategorię argumentów. \\
\hline
\end{tabular}
\end{center}

\subsection{Idiom \emph{pass by value, then move}}

\begin{lstlisting}
#include <string>
#include <utility>

class User {
    std::string name_;
public:
    explicit User(std::string name)
        : name_(std::move(name)) {}
};

std::string a = "Alice";
User first(a);                    // copy into parameter, then move
User second(std::string{"Bob"});  // parameter from prvalue, then move
\end{lstlisting}

To dobre rozwiązanie, gdy funkcja \emph{zawsze} potrzebuje własnej
kopii/przeniesienia argumentu. Nie jest jednak magicznie
najszybsze w każdym przypadku: dla lvalue może wymagać kopiowania
i jeszcze jednego przeniesienia, a duże obiekty niemożliwe do
przeniesienia mogą wymagać osobnego projektu API.

\subsection{Przekazywanie własności \texttt{unique\_ptr}}

\begin{lstlisting}
#include <memory>
#include <utility>

void adopt(std::unique_ptr<int> p) {
    // Ownership transferred into this function.
}

int main() {
    auto p = std::make_unique<int>(42);
    adopt(std::move(p));
    // p == nullptr after moving this unique_ptr
}
\end{lstlisting}

Przekazanie \texttt{unique\_ptr} przez wartość oznacza przejęcie
własności. Jeżeli funkcja ma jedynie korzystać z obiektu,
zwykle lepiej przyjąć \texttt{T\&} lub \texttt{const T\&}.
Temat własności zasobów omówiono w rozdziale~3.

\section{Rvalue-qualified methods i świadome udostępnianie zasobów}
\label{sec:move-ref-qualified}

C++ pozwala ograniczać metodę do lvalue albo rvalue odbiorcy,
dopisując \texttt{\&} lub \texttt{\&\&} po liście parametrów.
Można to wykorzystać np. do udostępnienia zawartości obiektu
bez kopiowania, gdy całe opakowanie jest tymczasowe.

\begin{lstlisting}
#include <string>
#include <utility>

class Message {
    std::string data_;
public:
    explicit Message(std::string s) : data_(std::move(s)) {}

    const std::string& text() const & noexcept {
        return data_;              // lvalue owner: borrow
    }
    std::string text() && noexcept {
        return std::move(data_);   // rvalue owner: consume
    }
};

Message m{"hello"};
const std::string& borrowed = m.text();
std::string owned = Message{"world"}.text();
\end{lstlisting}

Wariant \texttt{\&\&} może bezpiecznie oddać wartość przez
przeniesienie. Gdybyśmy zwracali w nim referencję do składowej
tymczasowego obiektu, powstałaby referencja wisząca.
Projektując takie API, warto rozważyć również wywołania na
\texttt{const} r-wartościach i ich wymagane zachowanie.

\section{Kontenery, iteratory przenoszące i wzorce użytkowe}

\subsection{Dodawanie elementów do \texttt{std::vector}}

\begin{lstlisting}
#include <string>
#include <utility>
#include <vector>

std::vector<std::string> words;
std::string word = "one";

words.push_back(word);             // copy from lvalue
words.push_back(std::move(word));  // move from lvalue cast to xvalue
words.emplace_back(3, 'x');       // construct "xxx" in place
\end{lstlisting}

Metoda \texttt{emplace\_back} konstruuje nowy element z podanych
argumentów. Nie musi być szybsza od \texttt{push\_back}.
Gotowy obiekt można przekazać z pomocą \texttt{std::move}.
Podczas realokacji kontener nadal może kopiować lub przenosić
wcześniej dodane elementy.

\subsection{Przenoszenie całego zakresu}

\begin{lstlisting}
#include <algorithm>
#include <iterator>
#include <string>
#include <vector>

std::vector<std::string> source{"a", "b"};
std::vector<std::string> dest;

std::move(source.begin(), source.end(),
          std::back_inserter(dest)); // algorithm std::move

// source still contains two valid strings, but their
// contents are not assumed.
\end{lstlisting}

Nazwa \texttt{std::move} występuje również dla algorytmu z nagłówka
\texttt{<algorithm>}. W przeciwieństwie do pojedynczego
\texttt{std::move(x)} z nagłówka \texttt{<utility>}, algorytm
przypisuje kolejno elementy z zakresu docelowego, korzystając
z możliwości ich przenoszenia.

\section{Typowe błędy i pułapki}
\label{sec:move-errors}

\subsection{Błąd 1: przekonanie, że \texttt{std::move} niszczy obiekt}

\texttt{std::move} jest rzutowaniem, a nie destruktorem. Obiekt
źródłowy nadal istnieje aż do końca własnego czasu życia.
Może zmienić stan, jeśli na jego podstawie wykonano
rzeczywistą operację przeniesienia.

\subsection{Błąd 2: \texttt{std::move} z argumentu przekazanego jako \texttt{const T\&}}

\begin{lstlisting}
#include <string>
#include <utility>

std::string ineffective(const std::string& x) {
    return std::move(x); // copies for std::string
}
\end{lstlisting}

Jeżeli funkcja potrzebuje zwrócić niezależną kopię, czytelniej
napisać \texttt{return x;}. Jeśli ma konsumować argument,
powinna otrzymać wartość lub odpowiednią modyfikowalną
referencję. Nie wolno usuwać \texttt{const} przez rzutowanie,
by wymusić przeniesienie z oryginalnie stałego obiektu.

\subsection{Błąd 3: użycie nazwanego \texttt{T\&\&} bez \texttt{std::move}}

\begin{lstlisting}
#include <memory>
#include <utility>

void take(std::unique_ptr<int>&& p) {
    // auto q = p;            // error: p is an lvalue
    auto q = std::move(p);    // explicit ownership transfer
}
\end{lstlisting}

\subsection{Błąd 4: przechowywanie referencji do argumentu tymczasowego}

\begin{lstlisting}
#include <string>
#include <string_view>

std::string_view dangerous() {
    return std::string{"temporary"}; // dangling view
}
\end{lstlisting}

\texttt{std::string\_view} jest widokiem, a nie właścicielem
pamięci. Semantyka przenoszenia nie rozwiązuje problemów
związanych z czasem życia. Gdy wynik funkcji ma być własnością
wywołującego, zwracaj \texttt{std::string} przez wartość.

\subsection{Błąd 5: przenoszenie nie gwarantuje zachowania wskaźników}

Przeniesienie kontenera może spowodować unieważnienie wcześniej
zachowanych iteratorów, referencji lub wskaźników zależnie od
typu kontenera, operacji i alokatora. Nie zakładaj ogólnej
zasady, że po \texttt{std::move} adresy danych zawsze pozostają
takie same. Sprawdź dokumentację odpowiedniego typu.

\subsection{Błąd 6: nadmierne stosowanie \texttt{std::forward}}

\texttt{std::forward<T>} ma sens głównie tam, gdzie \texttt{T}
pochodzi z dedukcji dla referencji przekazującej. Przy zwykłej
lokalnej zmiennej przeznaczonej do przekazania własności
należy używać \texttt{std::move}; przy zwracaniu lokalnej
zmiennej przez wartość zazwyczaj żadnego z nich.

\section{Przykład przekrojowy: kolejka zadań}
\label{sec:move-project}

Rozważmy prosty bufor zadań przechowujący funkcje do
późniejszego wykonania. Użyjemy \texttt{std::unique\_ptr},
aby zadanie jednoznacznie posiadało zasób. Poniższy przykład
jest niezależny od wielowątkowości; \emph{kolejka zadań}
nie uruchamia niczego w osobnym wątku.

\begin{lstlisting}
#include <functional>
#include <iostream>
#include <memory>
#include <utility>
#include <vector>

class Tasks {
    std::vector<std::function<void()>> items_;
public:
    template<class F>
    void add(F&& f) {
        items_.emplace_back(std::forward<F>(f));
    }

    void run() {
        for (auto& f : items_) f();
        items_.clear();
    }
};

int main() {
    Tasks tasks;
    auto value = std::make_shared<int>(10);
    tasks.add([value] { std::cout << *value << '\n'; });
    tasks.run();
}
\end{lstlisting}

\texttt{std::function<void()>} w C++20 wymaga, aby przechowywany
obiekt wywoływalny był kopiowalny. Dlatego w przykładzie lambda
przechowuje \texttt{shared\_ptr}. Gdybyśmy zamiast niego
przechwycili przez przeniesienie \texttt{unique\_ptr}, lambda
stałaby się niekopiowalna i taki obiekt nie spełniłby
wymagań \texttt{std::function} w C++20. C++23 wprowadza
\texttt{std::move\_only\_function}, która rozwiązuje ten problem.
To pokazuje, że prawidłowy wybór sposobu przenoszenia zależy
od kontraktów pozostałych typów.

\section{Praktyczna ściąga i zalecenia projektowe}

\begin{center}
\begin{tabular}{|p{4.0cm}|p{8.4cm}|}
\hline
\textbf{Sytuacja} & \textbf{Zalecenie} \\
\hline
Zwracam utworzony obiekt & Zwróć go przez wartość; korzystaj z elizji lub NRVO. \\
\hline
Przenoszę zasób z lokalnego obiektu & Użyj \texttt{std::move(x)} w miejscu przekazania własności. \\
\hline
Piszę wrapper o argumentach \texttt{Args\&\&...} & Wykorzystaj \texttt{std::forward<Args>(args)...}. \\
\hline
Klasa ma \texttt{vector}/\texttt{string} jako składowe & Preferuj regułę zera i generowane operacje. \\
\hline
Klasa posiada goły wskaźnik do alokacji & Preferuj zmianę projektu na RAII; jeśli zarządzasz nim ręcznie, rozważ wszystkie operacje specjalne. \\
\hline
Kontener przenosi elementy przy realokacji & Właściwy, bezpieczny \texttt{noexcept} w konstruktorze move pomaga w gwarancjach wyjątkowych. \\
\hline
Argument będzie tylko odczytywany & Zwykle \texttt{const T\&} dla dużego \texttt{T}, a małe typy przez wartość. \\
\hline
Funkcja przejmuje \texttt{unique\_ptr} & Przyjmij go przez wartość; wywołujący świadomie używa \texttt{std::move}. \\
\hline
\end{tabular}
\end{center}

\paragraph{Najważniejsze rozróżnienie.}
\begin{enumerate}
\item \texttt{std::move} wyraża zgodę na potraktowanie obiektu jako źródła przeniesienia, ale samo zasobów nie przenosi.
\item \texttt{std::forward<T>} zachowuje kategorię argumentu dedukowanego w szablonie.
\item Konstruktor przenoszący to zwykle operacja \texttt{T(T\&\&)}, ale może wykonywać dowolną operację zgodną z kontraktem typu.
\item \texttt{T\&\&} jako typ parametru nie zawsze jest referencją przekazującą.
\item Nazwany parametr typu \texttt{T\&\&} jest wyrażeniem \texttt{lvalue}.
\item Zwracanie przez wartość jest często najlepszym sposobem przekazywania wyniku.
\item Kiedy tylko się da, stosuj regułę zera zamiast własnych pięciu funkcji specjalnych.
\end{enumerate}

\paragraph{Pytania kontrolne.}
\begin{enumerate}
\item Dlaczego wywołanie \texttt{g(x)} w funkcji \texttt{f(T\&\& x)} traktuje \texttt{x} jako lvalue?
\item Kiedy \texttt{std::move(x)} kończy się kopiowaniem zamiast przeniesienia?
\item Czym różnią się \texttt{return T\{\};} i \texttt{return local;} w C++20?
\item Dlaczego przy powiększaniu \texttt{vector} konstruktor move z \texttt{noexcept} ma znaczenie?
\item Kiedy użyć \texttt{std::forward<T>(x)}, a kiedy \texttt{std::move(x)}?
\item Dlaczego metoda zwracająca \texttt{const std::string\&} nie powinna ujawniać referencji do składowej obiektu tymczasowego?
\end{enumerate}

\paragraph{Dalsza lektura.} \emph{cppreference}: \textit{Value categories},
\textit{Move constructors}, \textit{Reference declaration},
\textit{Copy elision}, \textit{std::move}, \textit{std::forward}.


\chapter[Semantyka przenoszenia w C++]{Semantyka przenoszenia, kategorie wartości i przekazywanie doskonałe}
\label{chap:move-semantics-cpp}

\section{Dlaczego wprowadzono semantykę przenoszenia?}

Rozważmy kontener przechowujący milion liczb. Skopiowanie go oznacza
zwykle przydzielenie nowej pamięci i skopiowanie każdego elementu.
Jeżeli dotychczasowy obiekt nie będzie już potrzebny, bardziej naturalna
jest operacja \emph{przekazania własności} jego zasobu zamiast kopiowania.
Od C++11 język udostępnia \emph{semantykę przenoszenia} (\textit{move
semantics}): programista może rozróżnić sytuację, w której obiekt źródłowy
ma pozostać niezmieniony, od sytuacji, w której jego zasoby wolno wykorzystać
ponownie.

\begin{lstlisting}
#include <string>
#include <utility>
#include <vector>

std::vector<std::string> a(100'000, "data");
auto b = a;                // copy: a remains intact

auto c = std::move(a);     // move may transfer a's allocation
// c now owns its elements; a is still valid, but its state
// after this move should not be assumed without a guarantee.
\end{lstlisting}

Nie oznacza to, że każde przeniesienie jest szybsze niż kopiowanie.
Dla typu \texttt{int} oba działania są w praktyce tożsame, a niektóre typy
kopiują część małych danych również podczas przenoszenia.
Kluczowa idea jest semantyczna: \emph{przeniesienie pozwala zmienić stan
źródła, a kopiowanie nie powinno go zmieniać}.

\paragraph{Zakres rozdziału.}
Zakładamy \textbf{C++20}. Kategorie wartości istnieją od początku języka,
referencje do r-wartości i \texttt{std::move} wprowadzono w C++11,
\emph{guaranteed copy elision} w C++17, a wiele przykładów poniżej
wykorzystuje koncepcje omówione w rozdziale~5.
Wszystkie przykłady zawierają komentarze angielskie, aby łatwo
przenosić je do projektów.

\section{Typ wyrażenia a kategoria wartości}
\label{sec:move-categories}

\textbf{Typ} i \textbf{kategoria wartości} są dwiema różnymi własnościami
wyrażenia. Typ mówi, jaki rodzaj danych otrzymujemy; kategoria mówi
m.in., czy wyrażenie identyfikuje istniejący obiekt i jak można
wiązać je z referencjami. Dokładnie trzy podstawowe kategorie to:
\emph{lvalue}, \emph{xvalue} i \emph{prvalue}.

\begin{center}
\begin{tabular}{|l|p{10.5cm}|}
\hline
\textbf{Kategoria} & \textbf{Intuicja i przykłady} \\
\hline
\texttt{lvalue} & Wyrażenie identyfikujące obiekt lub funkcję; np. \texttt{x}, \texttt{*ptr}, \texttt{++x} dla wbudowanego \texttt{int}. \\
\hline
\texttt{xvalue} & \emph{Expiring value}: wyrażenie identyfikujące obiekt, którego zasoby mogą być ponownie wykorzystane; np. \texttt{std::move(x)}. \\
\hline
\texttt{prvalue} & \emph{Pure rvalue}: wyrażenie obliczające wartość lub inicjalizujące obiekt; np. \texttt{42}, \texttt{x + 1}, \texttt{std::string("hi")}. \\
\hline
\end{tabular}
\end{center}

Pozostałe dwie nazwy oznaczają sumy kategorii:
\[
\underbrace{\mathrm{glvalue}}_{\mathrm{lvalue}\,\cup\,\mathrm{xvalue}},
\qquad
\underbrace{\mathrm{rvalue}}_{\mathrm{prvalue}\,\cup\,\mathrm{xvalue}}.
\]
Nie należy utożsamiać \texttt{rvalue} z tymczasowym obiektem ani
\texttt{lvalue} z wyrażeniem znajdującym się po lewej stronie przypisania.
Są to techniczne kategorie zdefiniowane przez język.

\subsection{Kilka wyrażeń o zaskakujących kategoriach}

\begin{lstlisting}
#include <string>
#include <utility>

int x = 7;
int* p = &x;

// x                 : lvalue
// *p                : lvalue
// (x)               : lvalue
// 7                 : prvalue
// x + 1             : prvalue
// std::move(x)      : xvalue
// std::string{"abc"}: prvalue

int& lref = x;
int&& rref = 9;

// lref              : lvalue
// rref              : lvalue (!) - it has a name
// std::move(rref)   : xvalue
\end{lstlisting}

To fundamentalna zasada: \textbf{nazwa zmiennej jest wyrażeniem lvalue,
nawet jeśli typ tej zmiennej to referencja rvalue \texttt{T\&\&}}.
Samo zadeklarowanie \texttt{T\&\&} nie oznacza, że każde użycie takiej
zmiennej będzie r-wartością.

\subsection{Badanie kategorii za pomocą \texttt{decltype((expr))}}

Dla prostego, nieujętego w nawiasy identyfikatora składnia
\texttt{decltype(x)} zwraca deklarowany typ. Dla ogólnego wyrażenia
\texttt{decltype((expr))} oddaje jego kategorię poprzez typ:
\[
\mathrm{decltype}((e))=
\left\{
\begin{array}{ll}
T\&, & \mbox{gdy } e \mbox{ jest lvalue typu } T,\\
T\&\&, & \mbox{gdy } e \mbox{ jest xvalue typu } T,\\
T, & \mbox{gdy } e \mbox{ jest prvalue typu } T.
\end{array}
\right.
\]

\begin{lstlisting}
#include <type_traits>
#include <utility>

int x = 10;
static_assert(std::is_same_v<decltype(x), int>);
static_assert(std::is_same_v<decltype((x)), int&>);
static_assert(std::is_same_v<decltype((std::move(x))), int&&>);
static_assert(std::is_same_v<decltype((x + 1)), int>);
\end{lstlisting}

\paragraph{Uwaga.} \texttt{decltype} stosuje szczególne reguły dla
nieujętych w nawiasy nazw i dostępu do składowych. Dlatego
\texttt{decltype(x)} oraz \texttt{decltype((x))} mogą zwrócić
różne typy. Nie jest to sprzeczność, lecz zamierzona cecha języka.

\subsection{Różnica między \texttt{auto} i \texttt{decltype(auto)}}

Zwykłe \texttt{auto} w deklaracji zmiennej najczęściej dedukuje typ
przez wartość i pomija referencję wynikającą z inicjalizatora.
Natomiast \texttt{decltype(auto)} stosuje reguły \texttt{decltype}
i może zachować kategorię wartości wyrażenia.

\begin{lstlisting}
#include <type_traits>

int n = 5;
auto copy = (n);           // int: independent value
decltype(auto) alias = (n); // int&: refers to n

static_assert(std::is_same_v<decltype(copy), int>);
static_assert(std::is_same_v<decltype(alias), int&>);

auto&& first = n;          // deduces int&
auto&& second = 5;         // deduces int&&
static_assert(std::is_same_v<decltype(first), int&>);
static_assert(std::is_same_v<decltype(second), int&&>);
\end{lstlisting}

W typie wyniku funkcji zapis \texttt{decltype(auto)} wymaga
ostrożności. Instrukcja \texttt{return (x);} może zwrócić
referencję do zmiennej lokalnej, która przestanie istnieć.
Zwykłe \texttt{auto} jest często bezpieczniejszym wyborem,
gdy funkcja ma zwrócić własny obiekt przez wartość.

\section{Referencje lvalue i rvalue}
\label{sec:move-references}

\begin{lstlisting}
#include <string>

std::string name = "Ada";
std::string& a = name;            // binds to mutable lvalue
const std::string& b = name;      // binds to lvalue
const std::string& c = "temporary"; // can bind to temporary
std::string&& d = std::string{"x"}; // binds to rvalue

// std::string& e = std::string{"x"}; // error
// std::string&& f = name;           // error
\end{lstlisting}

Referencja \texttt{T\&} umożliwia pracę z istniejącym, zwykle
nazwanym obiektem; \texttt{T\&\&} może wiązać się z r-wartością.
\texttt{const T\&} wiąże się zarówno z lvalue, jak i z rvalue,
nie dając prawa do modyfikowania obiektu poprzez tę referencję.

\subsection{Wybór przeciążenia}

\begin{lstlisting}
#include <iostream>
#include <string>
#include <utility>

void inspect(const std::string&) {
    std::cout << "const lvalue reference\n";
}
void inspect(std::string&&) {
    std::cout << "rvalue reference\n";
}

int main() {
    std::string text = "hello";
    inspect(text);             // const lvalue reference
    inspect(std::string{"x"}); // rvalue reference
    inspect(std::move(text));  // rvalue reference
}
\end{lstlisting}

Istnienie przeciążenia \texttt{f(T\&\&)} samo w sobie nie przesądza
o rzeczywistym przeniesieniu czegokolwiek. Funkcja może tylko
odczytać argument; \emph{move} jest operacją wykonywaną przez
konstruktor lub przypisanie, a nie przez mechanizm wiązania referencji.

\subsection{Wydłużenie życia obiektu tymczasowego}

W typowych deklaracjach lokalnych związanie obiektu tymczasowego
bezpośrednio z referencją może wydłużyć jego życie do końca życia
tej referencji:
\begin{lstlisting}
#include <string>

const std::string& a = std::string{"safe until a ends"};
std::string&& b = std::string{"safe until b ends"};
\end{lstlisting}
Nie wolno jednak uogólniać tej reguły na zwracanie referencji ani
na dowolne przekazywanie referencji dalej. Przykładowo:
\begin{lstlisting}
#include <string>

const std::string& broken() {
    return std::string{"dangling"}; // dangling reference
}
\end{lstlisting}
Zwrócona referencja nie utrzymuje tego obiektu przy życiu.
Jeżeli funkcja ma \emph{wytworzyć} wynik, zwykle powinna zwracać go
przez wartość.

\section{Co naprawdę robi \texttt{std::move}?}
\label{sec:move-std-move}

Nazwa jest nieco myląca: \texttt{std::move(x)} \textbf{nie przenosi
zasobów samodzielnie}. W przybliżeniu wykonuje rzutowanie do
referencji rvalue odpowiedniego typu. Dopiero dalszy konstruktor,
operator przypisania lub inna funkcja może potraktować wynik
jako obiekt nadający się do przeniesienia.

\begin{lstlisting}
#include <type_traits>
#include <utility>

template<class T>
constexpr std::remove_reference_t<T>&&
my_move(T&& value) noexcept {
    return static_cast<std::remove_reference_t<T>&&>(value);
}
\end{lstlisting}

Powyższa implementacja ilustruje istotę działania
\texttt{std::move}. W rzeczywistym kodzie korzystamy z funkcji
standardowej z nagłówka \texttt{<utility>}.

\begin{lstlisting}
#include <string>
#include <utility>

std::string first = "long text";
std::string second = std::move(first);
// A move constructor is selected for second.
// first is valid; do not assume its exact contents.

first = "reused";             // valid: assignment
first.clear();                // valid: clear()
\end{lstlisting}

\paragraph{Stan obiektu po przeniesieniu.}
Typy biblioteczne zwykle gwarantują stan \emph{poprawny, lecz
nieokreślony co do zawartości} (\textit{valid but unspecified}), o ile
kontrakt konkretnego typu nie obiecuje więcej. Takiego obiektu
można bezpiecznie zniszczyć, przypisać mu nową wartość albo
wykonać operację, której wymagania wstępne na to pozwalają.
Nie należy jednak \emph{zakładać}, że każdy przeniesiony
\texttt{std::string} jest pusty. Dla niektórych typów, np.
\texttt{std::unique\_ptr}, stan po przeniesieniu określono dokładniej.

\subsection{Problem obiektu \texttt{const}}

\begin{lstlisting}
#include <string>
#include <utility>

const std::string original = "cannot be emptied here";
std::string other = std::move(original);
// Typically: copy, not move (const string&& cannot bind to
// string's non-const move constructor).
\end{lstlisting}

\texttt{std::move(original)} ma w tym przypadku typ
\texttt{const std::string\&\&}. Typowy konstruktor przenoszący
przyjmuje \texttt{std::string\&\&}, więc nie może otrzymać
argumentu \texttt{const}. Kompilator wybiera konstruktor kopiujący
przyjmujący \texttt{const std::string\&}. \textbf{Samo wystąpienie
\texttt{std::move} w kodzie nie gwarantuje przeniesienia}.\footnote{Dla
typów niekopiowalnych ten sam zapis może skończyć się błędem
kompilacji.}

\section{Konstruktor przenoszący i przenoszący operator przypisania}
\label{sec:move-special-members}

W klasie \texttt{T} konstruktor przenoszący ma zwykle sygnaturę
\texttt{T(T\&\&)}, natomiast przenoszący operator przypisania:
\texttt{T\& operator=(T\&\&)}. Nie należy mylić ich z kopiowaniem:
\begin{lstlisting}
struct Item {
    Item() = default;
    Item(const Item&);            // copy constructor
    Item& operator=(const Item&); // copy assignment
    Item(Item&&) noexcept;        // move constructor
    Item& operator=(Item&&) noexcept; // move assignment
    ~Item();                      // destructor
};
\end{lstlisting}

\subsection{Przykład z obserwacją operacji}
\label{sec:move-trace}

\begin{lstlisting}
#include <iostream>
#include <utility>

struct Trace {
    Trace() { std::cout << "default\n"; }
    Trace(const Trace&) { std::cout << "copy ctor\n"; }
    Trace(Trace&&) noexcept { std::cout << "move ctor\n"; }
    Trace& operator=(const Trace&) {
        std::cout << "copy assign\n";
        return *this;
    }
    Trace& operator=(Trace&&) noexcept {
        std::cout << "move assign\n";
        return *this;
    }
};

int main() {
    Trace a;                    // default
    Trace b = a;                // copy ctor
    Trace c = std::move(a);     // move ctor
    b = c;                      // copy assign
    b = std::move(c);           // move assign
}
\end{lstlisting}

Konstruktor tworzy \emph{nowy} obiekt, a operator przypisania
zmienia już istniejący. W tym przykładzie \texttt{Trace} jedynie
wypisuje komunikaty. W klasie zarządzającej zasobem konstruktor
przenoszący powinien przekazać własność zasobu, a operator
przypisania musi dodatkowo poprawnie pozbyć się zasobu
posiadanego wcześniej przez cel.

\subsection{Pełny przykład: klasa samodzielnie zarządzająca tablicą}
\label{sec:move-raw-buffer}

To przykład \emph{edukacyjny}. W rzeczywistym projekcie zwykle
lepiej użyć standardowego kontenera, np. \texttt{std::vector<int>}.

\begin{lstlisting}
#include <algorithm>
#include <cstddef>
#include <utility>

class RawBuffer {
    std::size_t size_ = 0;
    int* data_ = nullptr;

public:
    explicit RawBuffer(std::size_t n)
        : size_(n), data_(n ? new int[n]{} : nullptr) {}

    ~RawBuffer() { delete[] data_; }

    RawBuffer(const RawBuffer& other)
        : RawBuffer(other.size_) {
        if (size_ != 0)
            std::copy_n(other.data_, size_, data_);
    }

    RawBuffer& operator=(const RawBuffer& other) {
        if (this != &other) {
            RawBuffer copy(other); // may throw; *this unchanged
            swap(copy);
        }
        return *this;
    }

    RawBuffer(RawBuffer&& other) noexcept
        : size_(std::exchange(other.size_, 0)),
          data_(std::exchange(other.data_, nullptr)) {}

    RawBuffer& operator=(RawBuffer&& other) noexcept {
        if (this != &other) {
            delete[] data_;
            size_ = std::exchange(other.size_, 0);
            data_ = std::exchange(other.data_, nullptr);
        }
        return *this;
    }

    void swap(RawBuffer& other) noexcept {
        using std::swap;
        swap(size_, other.size_);
        swap(data_, other.data_);
    }

    std::size_t size() const noexcept { return size_; }
    int& operator[](std::size_t i) { return data_[i]; }
    const int& operator[](std::size_t i) const { return data_[i]; }
};
\end{lstlisting}

\texttt{std::exchange(obj, value)} zwraca starą wartość \texttt{obj}
i zastępuje ją wartością \texttt{value}. Konstruktor przenoszący
przejmuje więc wskaźnik i rozmiar, a źródło ustawia w jawnie
określony stan pusty. Operator przenoszącego przypisania
usuwa dotychczasową tablicę celu, przejmuje nową i zabezpiecza
przypadek \texttt{b = std::move(b)}. W implementacji kopiowania
widać z kolei \emph{silną gwarancję wyjątkową}: jeśli utworzenie
kopii nie powiedzie się, pierwotny obiekt pozostaje bez zmian.

\section{Reguła pięciu, reguła trzech i reguła zera}
\label{sec:move-rule-of-five}

Klasy zarządzające zasobem często muszą kontrolować pięć operacji:
(1) destruktor, (2) konstruktor kopiujący, (3) kopiujący operator
przypisania, (4) konstruktor przenoszący oraz (5) przenoszący
operator przypisania. To \emph{reguła pięciu}.

Historyczna \emph{reguła trzech} odnosiła się do destruktora,
konstruktora kopiującego i kopiującego operatora przypisania,
zanim język udostępnił move semantics.

\paragraph{Ważne: deklarowanie destruktora może blokować generowanie
konstruktora przenoszącego.}

\begin{lstlisting}
#include <memory>
#include <type_traits>

struct WithDestructor {
    std::unique_ptr<int> p;
    ~WithDestructor() = default;
};

static_assert(!std::is_move_constructible_v<WithDestructor>);

struct WithMove {
    std::unique_ptr<int> p;
    ~WithMove() = default;
    WithMove() = default;
    WithMove(WithMove&&) noexcept = default;
    WithMove& operator=(WithMove&&) noexcept = default;
};
static_assert(std::is_move_constructible_v<WithMove>);
\end{lstlisting}

W \texttt{WithDestructor} użytkownik zadeklarował destruktor,
więc kompilator nie deklaruje automatycznie operacji przenoszących.
Kopiowanie jest niedostępne z powodu \texttt{unique\_ptr}.
Samo \texttt{= default} przy destruktorze nie usuwa tej reguły.

\subsection{Najlepsza praktyka: reguła zera}

\emph{Reguła zera} mówi, że klasy, które nie są bezpośrednio
odpowiedzialne za posiadanie surowych zasobów, zwykle nie powinny
same definiować żadnej z pięciu operacji.

\begin{lstlisting}
#include <string>
#include <vector>

struct Dataset {
    std::string title;
    std::vector<double> values;
};

Dataset original{"results", {1.0, 2.0, 3.0}};
Dataset copy = original;           // deep copy of vector data
Dataset moved = std::move(original); // member-wise move
\end{lstlisting}

To kontenery standardowe odpowiadają za przydział i zwalnianie
pamięci. Automatycznie generowane operacje \texttt{Dataset}
wykonują odpowiednie operacje na jego składowych. \textbf{Zalecenie:}
zacznij od reguły zera; własny konstruktor przenoszący pisz
jedynie wtedy, gdy zachowanie domyślne nie jest właściwe.

\section{Przenoszenie a wyjątki: \texttt{noexcept}}
\label{sec:move-noexcept}

Kontenery takie jak \texttt{std::vector<T>} przy realokacji muszą
przenieść lub skopiować dotychczasowe elementy do nowego obszaru.
Jeśli przenoszenie \texttt{T} może rzucać wyjątki, a \texttt{T}
można kopiować, biblioteka może wybrać kopiowanie, by zachować
silniejsze gwarancje wyjątkowe. Właśnie dlatego poprawne operacje
przenoszące, które nie mogą rzucać wyjątków, warto deklarować
\texttt{noexcept}.

\begin{lstlisting}
#include <type_traits>
#include <utility>

struct Resource {
    Resource() = default;
    Resource(const Resource&) = default;
    Resource(Resource&&) noexcept = default;
};

static_assert(std::is_nothrow_move_constructible_v<Resource>);

// move_if_noexcept can provide an lvalue reference if
// moving might throw and a copy is available.
\end{lstlisting}

\texttt{std::move\_if\_noexcept(x)} zwraca odpowiednio
\texttt{T\&\&} albo \texttt{const T\&}, zależnie od własności
konstruktorów danego typu. To narzędzie używane przy implementowaniu
kontenerów, a nie uniwersalny zamiennik \texttt{std::move}.

\paragraph{Ostrożnie z \texttt{noexcept}.}
Wyjście wyjątku poza funkcję zadeklarowaną jako
\texttt{noexcept} wywołuje \texttt{std::terminate}.
Dlatego bezwyjątkowość musi wynikać z rzeczywistych operacji składowych.

\section{Zwijanie referencji (\textit{reference collapsing})}
\label{sec:move-ref-collapse}

Gdy referencje występują pośrednio poprzez aliasy lub
podstawienie argumentów szablonu, stosuje się \emph{zwijanie
referencji}. Wszystkie przypadki przedstawia tabela:

\begin{center}
\begin{tabular}{|c|c|c|}
\hline
\textbf{Typ pośredni} & \textbf{Zapis} & \textbf{Wynik} \\
\hline
\texttt{T\&} i \texttt{\&} & \texttt{T\& \&} & \texttt{T\&} \\
\texttt{T\&} i \texttt{\&\&} & \texttt{T\& \&\&} & \texttt{T\&} \\
\texttt{T\&\&} i \texttt{\&} & \texttt{T\&\& \&} & \texttt{T\&} \\
\texttt{T\&\&} i \texttt{\&\&} & \texttt{T\&\& \&\&} & \texttt{T\&\&} \\
\hline
\end{tabular}
\end{center}

Krótko: jeżeli \emph{choć jeden} składnik jest \texttt{\&},
wynik jest \texttt{\&}; tylko \texttt{\&\&} połączone
z \texttt{\&\&} daje \texttt{\&\&}.

\begin{lstlisting}
#include <type_traits>

using L = int&;
using R = int&&;

static_assert(std::is_same_v<L&, int&>);
static_assert(std::is_same_v<L&&, int&>);
static_assert(std::is_same_v<R&, int&>);
static_assert(std::is_same_v<R&&, int&&>);
\end{lstlisting}

Zwijanie referencji jest podstawą \emph{perfect forwarding}:
dzięki niemu jeden szablon może przyjmować zarówno lvalue,
jak i rvalue i przekazywać tę różnicę dalej.

\section{Referencje przekazujące i dedukcja \texttt{T\&\&}}
\label{sec:move-forwarding-ref}

Zapis \texttt{T\&\&} ma dwa znaczenia w zależności od kontekstu.
Jeśli \texttt{T} jest dedukowanym, niekwalifikowanym parametrem
typu \emph{tego samego} szablonu funkcji, \texttt{T\&\&}
to \emph{referencja przekazująca} (\textit{forwarding reference}).
W pozostałych typowych kontekstach jest to po prostu
referencja rvalue.

\begin{lstlisting}
#include <string>

template<class T>
void forwarding(T&& x);         // forwarding reference

template<class T>
void not_forwarding(const T&& x); // rvalue reference

void fixed(std::string&& x);    // rvalue reference

template<class T>
struct Holder {
    void receive(T&& x);        // NOT forwarding: T comes
                               // from the class template

    template<class U>
    void relay(U&& x);          // forwarding: U is deduced
};
\end{lstlisting}

Dla wywołania \texttt{forwarding(x)} z lvalue typu \texttt{int}
kompilator dedukuje \texttt{T = int\&}, zatem parametr ma po
zwinięciu typ \texttt{int\&}. Dla \texttt{forwarding(5)}
dedukuje \texttt{T = int}, a parametr ma typ \texttt{int\&\&}.

\begin{lstlisting}
#include <type_traits>

template<class T>
void check(T&& x) {
    if constexpr (std::is_lvalue_reference_v<T>) {
        // Called with an lvalue.
    } else {
        // Called with an rvalue.
    }
}
\end{lstlisting}

\paragraph{Drobne wyjątki.} \texttt{auto\&\&} zwykle również
pozwala na przekazywanie kategorii, lecz dedukcja \texttt{auto}
ma szczególne reguły dla list inicjalizacyjnych.
\texttt{const T\&\&} nie jest referencją przekazującą.

\section{Perfect forwarding i \texttt{std::forward}}
\label{sec:move-perfect-forwarding}

Przez \emph{perfect forwarding} rozumiemy przekazanie argumentu
do następnej funkcji przy zachowaniu kategorii wartości,
a także odpowiednich kwalifikatorów typu.

\begin{lstlisting}
#include <iostream>
#include <utility>

void target(int&)  { std::cout << "lvalue\n"; }
void target(int&&) { std::cout << "rvalue\n"; }

template<class T>
void incorrect(T&& value) {
    target(value); // value is a named variable: always lvalue
}

template<class T>
void correct(T&& value) {
    target(std::forward<T>(value));
}

int main() {
    int x = 1;
    incorrect(2); // lvalue (!)
    correct(x);   // lvalue
    correct(2);   // rvalue
}
\end{lstlisting}

W przybliżeniu \texttt{std::forward<T>(value)} zwraca
lvalue, jeżeli \texttt{T} wydedukowano jako referencję lvalue,
i xvalue w przeciwnym razie. \texttt{std::move(value)} zawsze
wymusza kategorię xvalue: w funkcji przekazującej byłoby to
błędem, ponieważ niepotrzebnie pozwalałoby na przeniesienie
zasobu z argumentu przekazanego jako lvalue.

\subsection{Przykład: fabryka obiektów}

\begin{lstlisting}
#include <memory>
#include <string>
#include <utility>

struct Person {
    std::string name;
    int age;

    Person(std::string n, int a)
        : name(std::move(n)), age(a) {}
};

template<class T, class... Args>
std::unique_ptr<T> make_object(Args&&... args) {
    return std::unique_ptr<T>(
        new T(std::forward<Args>(args)...)
    );
}

int main() {
    std::string first = "Ada";
    auto p = make_object<Person>(first, 37); // first copied
    auto q = make_object<Person>(std::move(first), 38);
}
\end{lstlisting}

To edukacyjna wersja fabryki. W produkcyjnym kodzie
\texttt{std::make\_unique<T>(...)} robi to samo w prostszy
sposób. W drugim wywołaniu użytkownik świadomie pozwolił na
przeniesienie \texttt{first}; w pierwszym pozostaje on
nienaruszony. \texttt{std::forward<Args>(args)...} rozwija pakiet
argumentów omówiony w rozdziale~4.

\subsection{Nie przekazuj tego samego zasobu dwa razy bez zastanowienia}

\begin{lstlisting}
#include <utility>

template<class T>
void bad_double_forward(T&& x) {
    consume(std::forward<T>(x));
    consume(std::forward<T>(x)); // may use moved-from object
}
\end{lstlisting}

Ten fragment jest \emph{schematem błędu}: zakładamy, że istnieje
funkcja \texttt{consume}. Jeśli pierwsze wywołanie wykorzysta
argument rvalue przez przeniesienie, drugie otrzyma już obiekt
po przeniesieniu. Referencję przekazującą traktuj zatem jako
przeznaczoną zwykle do \emph{jednego} przekazania dalej.

\section{Perfect forwarding a przeciążenia i szablony klas}

Konstruktor z referencją przekazującą potrafi być bardzo zachłanny
podczas rozstrzygania przeciążeń. Może np. wygrać z konstruktorem
kopiującym przy argumentach nie-\texttt{const}.

\begin{lstlisting}
#include <string>
#include <utility>

struct Named {
    std::string name;

    Named(const Named&) = default;

    template<class U>
    explicit Named(U&& u)
        : name(std::forward<U>(u)) {}
};

// Named a{"text"};
// Named b(a);  // may select the constructor template,
//              // then fail to initialize std::string from Named&
\end{lstlisting}

Zwykle lepiej zacząć od zwykłych przeciążeń albo konstruktora
przyjmującego konkretny typ. Jeżeli szablon jest potrzebny,
można ograniczyć go koncepcją:

\begin{lstlisting}
#include <concepts>
#include <string>
#include <type_traits>
#include <utility>

struct SafeNamed {
    std::string name;

    SafeNamed(const SafeNamed&) = default;
    SafeNamed(SafeNamed&&) = default;

    template<class U>
      requires (!std::same_as<std::remove_cvref_t<U>, SafeNamed>
                && std::constructible_from<std::string, U&&>)
    explicit SafeNamed(U&& u)
        : name(std::forward<U>(u)) {}
};
\end{lstlisting}

Takie powiązanie rozdziału o perfect forwarding z koncepcjami
z rozdziału~5 jest bardzo przydatne
w bibliotekach generycznych.

\section{Copy elision: kiedy nie ma ani kopii, ani przeniesienia?}
\label{sec:move-elision}

Prosty model \emph{zwrócenie przez wartość = kopia lub move}
jest niepoprawny w nowoczesnym C++. Od C++17 pewne wyrażenia
\texttt{prvalue} inicjalizują wynik \emph{bezpośrednio w miejscu
docelowym}. Nie istnieje dodatkowy obiekt, który trzeba by
skopiować lub przenieść.

\begin{lstlisting}
struct NonMovable {
    NonMovable() = default;
    NonMovable(const NonMovable&) = delete;
    NonMovable(NonMovable&&) = delete;
};

NonMovable create() {
    return NonMovable{}; // guaranteed elision, C++17+
}

int main() {
    NonMovable x = create(); // OK in C++17+
}
\end{lstlisting}

Nie jest to zwykła optymalizacja, którą kompilator \emph{może}
wyłączyć: w pokazanym przypadku wynika to z semantyki języka.

\subsection{NRVO a gwarantowana elizja}

\begin{lstlisting}
#include <string>

std::string first() {
    return std::string{"hello"}; // mandatory direct construction
}

std::string second() {
    std::string s = "hello";
    return s;                    // NRVO may occur
}
\end{lstlisting}

\textit{Named Return Value Optimization} (NRVO) dotyczy
zwracania nazwanej lokalnej zmiennej. W C++20 kompilator \emph{może}
pominąć utworzenie dodatkowego obiektu. Jeśli NRVO nie nastąpi,
wyrażenie \texttt{s} w instrukcji \texttt{return} może
uczestniczyć w niejawnej próbie przeniesienia.

\paragraph{Nie pisz automatycznie \texttt{return std::move(s);}.}

\begin{lstlisting}
std::string make_text() {
    std::string s = "result";
    return s;             // allows NRVO
}

std::string less_good() {
    std::string s = "result";
    return std::move(s);  // inhibits NRVO; move may be required
}
\end{lstlisting}

\texttt{std::move(s)} zmienia wynik w xvalue i tym samym
nie spełnia warunków NRVO dotyczących nazwy obiektu.
Jeśli funkcja zwraca własną lokalną zmienną przez wartość,
zazwyczaj najlepsze jest \texttt{return s;}.

\section{Argumenty funkcji: kiedy kopiować, kiedy przenosić?}
\label{sec:move-api}

Wybór sposobu przyjmowania argumentu jest częścią kontraktu API.
Nie chodzi jedynie o wydajność, ale przede wszystkim o własność
i intencję.

\begin{center}
\begin{tabular}{|p{3.4cm}|p{4.3cm}|p{4.7cm}|}
\hline
\textbf{Parametr} & \textbf{Znaczenie} & \textbf{Kiedy używać} \\
\hline
\texttt{const T\&} & Odczyt bez przejmowania własności. & Duże obiekty, gdy wystarczy tylko odczyt. \\
\hline
\texttt{T\&} & Modyfikacja obiektu należącego do wywołującego. & Parametry wyjściowe i operacje in-place. \\
\hline
\texttt{T} & Własna wartość; może być skopiowana albo przeniesiona. & Małe typy i funkcje mające zachować wartość. \\
\hline
\texttt{T\&\&} & Dostęp do argumentu będącego r-wartością. & Przeciążenia zależne od kategorii oraz operacje konsumujące. \\
\hline
\texttt{T\&\&} z dedukcją \texttt{T} & Referencja przekazująca. & Opakowania i fabryki zachowujące kategorię argumentów. \\
\hline
\end{tabular}
\end{center}

\subsection{Idiom \emph{pass by value, then move}}

\begin{lstlisting}
#include <string>
#include <utility>

class User {
    std::string name_;
public:
    explicit User(std::string name)
        : name_(std::move(name)) {}
};

std::string a = "Alice";
User first(a);                    // copy into parameter, then move
User second(std::string{"Bob"});  // parameter from prvalue, then move
\end{lstlisting}

To dobre rozwiązanie, gdy funkcja \emph{zawsze} potrzebuje własnej
kopii/przeniesienia argumentu. Nie jest jednak magicznie
najszybsze w każdym przypadku: dla lvalue może wymagać kopiowania
i jeszcze jednego przeniesienia, a duże obiekty niemożliwe do
przeniesienia mogą wymagać osobnego projektu API.

\subsection{Przekazywanie własności \texttt{unique\_ptr}}

\begin{lstlisting}
#include <memory>
#include <utility>

void adopt(std::unique_ptr<int> p) {
    // Ownership transferred into this function.
}

int main() {
    auto p = std::make_unique<int>(42);
    adopt(std::move(p));
    // p == nullptr after moving this unique_ptr
}
\end{lstlisting}

Przekazanie \texttt{unique\_ptr} przez wartość oznacza przejęcie
własności. Jeżeli funkcja ma jedynie korzystać z obiektu,
zwykle lepiej przyjąć \texttt{T\&} lub \texttt{const T\&}.
Temat własności zasobów omówiono w rozdziale~3.

\section{Rvalue-qualified methods i świadome udostępnianie zasobów}
\label{sec:move-ref-qualified}

C++ pozwala ograniczać metodę do lvalue albo rvalue odbiorcy,
dopisując \texttt{\&} lub \texttt{\&\&} po liście parametrów.
Można to wykorzystać np. do udostępnienia zawartości obiektu
bez kopiowania, gdy całe opakowanie jest tymczasowe.

\begin{lstlisting}
#include <string>
#include <utility>

class Message {
    std::string data_;
public:
    explicit Message(std::string s) : data_(std::move(s)) {}

    const std::string& text() const & noexcept {
        return data_;              // lvalue owner: borrow
    }
    std::string text() && noexcept {
        return std::move(data_);   // rvalue owner: consume
    }
};

Message m{"hello"};
const std::string& borrowed = m.text();
std::string owned = Message{"world"}.text();
\end{lstlisting}

Wariant \texttt{\&\&} może bezpiecznie oddać wartość przez
przeniesienie. Gdybyśmy zwracali w nim referencję do składowej
tymczasowego obiektu, powstałaby referencja wisząca.
Projektując takie API, warto rozważyć również wywołania na
\texttt{const} r-wartościach i ich wymagane zachowanie.

\section{Kontenery, iteratory przenoszące i wzorce użytkowe}

\subsection{Dodawanie elementów do \texttt{std::vector}}

\begin{lstlisting}
#include <string>
#include <utility>
#include <vector>

std::vector<std::string> words;
std::string word = "one";

words.push_back(word);             // copy from lvalue
words.push_back(std::move(word));  // move from lvalue cast to xvalue
words.emplace_back(3, 'x');       // construct "xxx" in place
\end{lstlisting}

Metoda \texttt{emplace\_back} konstruuje nowy element z podanych
argumentów. Nie musi być szybsza od \texttt{push\_back}.
Gotowy obiekt można przekazać z pomocą \texttt{std::move}.
Podczas realokacji kontener nadal może kopiować lub przenosić
wcześniej dodane elementy.

\subsection{Przenoszenie całego zakresu}

\begin{lstlisting}
#include <algorithm>
#include <iterator>
#include <string>
#include <vector>

std::vector<std::string> source{"a", "b"};
std::vector<std::string> dest;

std::move(source.begin(), source.end(),
          std::back_inserter(dest)); // algorithm std::move

// source still contains two valid strings, but their
// contents are not assumed.
\end{lstlisting}

Nazwa \texttt{std::move} występuje również dla algorytmu z nagłówka
\texttt{<algorithm>}. W przeciwieństwie do pojedynczego
\texttt{std::move(x)} z nagłówka \texttt{<utility>}, algorytm
przypisuje kolejno elementy z zakresu docelowego, korzystając
z możliwości ich przenoszenia.

\section{Typowe błędy i pułapki}
\label{sec:move-errors}

\subsection{Błąd 1: przekonanie, że \texttt{std::move} niszczy obiekt}

\texttt{std::move} jest rzutowaniem, a nie destruktorem. Obiekt
źródłowy nadal istnieje aż do końca własnego czasu życia.
Może zmienić stan, jeśli na jego podstawie wykonano
rzeczywistą operację przeniesienia.

\subsection{Błąd 2: \texttt{std::move} z argumentu przekazanego jako \texttt{const T\&}}

\begin{lstlisting}
#include <string>
#include <utility>

std::string ineffective(const std::string& x) {
    return std::move(x); // copies for std::string
}
\end{lstlisting}

Jeżeli funkcja potrzebuje zwrócić niezależną kopię, czytelniej
napisać \texttt{return x;}. Jeśli ma konsumować argument,
powinna otrzymać wartość lub odpowiednią modyfikowalną
referencję. Nie wolno usuwać \texttt{const} przez rzutowanie,
by wymusić przeniesienie z oryginalnie stałego obiektu.

\subsection{Błąd 3: użycie nazwanego \texttt{T\&\&} bez \texttt{std::move}}

\begin{lstlisting}
#include <memory>
#include <utility>

void take(std::unique_ptr<int>&& p) {
    // auto q = p;            // error: p is an lvalue
    auto q = std::move(p);    // explicit ownership transfer
}
\end{lstlisting}

\subsection{Błąd 4: przechowywanie referencji do argumentu tymczasowego}

\begin{lstlisting}
#include <string>
#include <string_view>

std::string_view dangerous() {
    return std::string{"temporary"}; // dangling view
}
\end{lstlisting}

\texttt{std::string\_view} jest widokiem, a nie właścicielem
pamięci. Semantyka przenoszenia nie rozwiązuje problemów
związanych z czasem życia. Gdy wynik funkcji ma być własnością
wywołującego, zwracaj \texttt{std::string} przez wartość.

\subsection{Błąd 5: przenoszenie nie gwarantuje zachowania wskaźników}

Przeniesienie kontenera może spowodować unieważnienie wcześniej
zachowanych iteratorów, referencji lub wskaźników zależnie od
typu kontenera, operacji i alokatora. Nie zakładaj ogólnej
zasady, że po \texttt{std::move} adresy danych zawsze pozostają
takie same. Sprawdź dokumentację odpowiedniego typu.

\subsection{Błąd 6: nadmierne stosowanie \texttt{std::forward}}

\texttt{std::forward<T>} ma sens głównie tam, gdzie \texttt{T}
pochodzi z dedukcji dla referencji przekazującej. Przy zwykłej
lokalnej zmiennej przeznaczonej do przekazania własności
należy używać \texttt{std::move}; przy zwracaniu lokalnej
zmiennej przez wartość zazwyczaj żadnego z nich.

\section{Przykład przekrojowy: kolejka zadań}
\label{sec:move-project}

Rozważmy prosty bufor zadań przechowujący funkcje do
późniejszego wykonania. Użyjemy \texttt{std::unique\_ptr},
aby zadanie jednoznacznie posiadało zasób. Poniższy przykład
jest niezależny od wielowątkowości; \emph{kolejka zadań}
nie uruchamia niczego w osobnym wątku.

\begin{lstlisting}
#include <functional>
#include <iostream>
#include <memory>
#include <utility>
#include <vector>

class Tasks {
    std::vector<std::function<void()>> items_;
public:
    template<class F>
    void add(F&& f) {
        items_.emplace_back(std::forward<F>(f));
    }

    void run() {
        for (auto& f : items_) f();
        items_.clear();
    }
};

int main() {
    Tasks tasks;
    auto value = std::make_shared<int>(10);
    tasks.add([value] { std::cout << *value << '\n'; });
    tasks.run();
}
\end{lstlisting}

\texttt{std::function<void()>} w C++20 wymaga, aby przechowywany
obiekt wywoływalny był kopiowalny. Dlatego w przykładzie lambda
przechowuje \texttt{shared\_ptr}. Gdybyśmy zamiast niego
przechwycili przez przeniesienie \texttt{unique\_ptr}, lambda
stałaby się niekopiowalna i taki obiekt nie spełniłby
wymagań \texttt{std::function} w C++20. C++23 wprowadza
\texttt{std::move\_only\_function}, która rozwiązuje ten problem.
To pokazuje, że prawidłowy wybór sposobu przenoszenia zależy
od kontraktów pozostałych typów.

\section{Praktyczna ściąga i zalecenia projektowe}

\begin{center}
\begin{tabular}{|p{4.0cm}|p{8.4cm}|}
\hline
\textbf{Sytuacja} & \textbf{Zalecenie} \\
\hline
Zwracam utworzony obiekt & Zwróć go przez wartość; korzystaj z elizji lub NRVO. \\
\hline
Przenoszę zasób z lokalnego obiektu & Użyj \texttt{std::move(x)} w miejscu przekazania własności. \\
\hline
Piszę wrapper o argumentach \texttt{Args\&\&...} & Wykorzystaj \texttt{std::forward<Args>(args)...}. \\
\hline
Klasa ma \texttt{vector}/\texttt{string} jako składowe & Preferuj regułę zera i generowane operacje. \\
\hline
Klasa posiada goły wskaźnik do alokacji & Preferuj zmianę projektu na RAII; jeśli zarządzasz nim ręcznie, rozważ wszystkie operacje specjalne. \\
\hline
Kontener przenosi elementy przy realokacji & Właściwy, bezpieczny \texttt{noexcept} w konstruktorze move pomaga w gwarancjach wyjątkowych. \\
\hline
Argument będzie tylko odczytywany & Zwykle \texttt{const T\&} dla dużego \texttt{T}, a małe typy przez wartość. \\
\hline
Funkcja przejmuje \texttt{unique\_ptr} & Przyjmij go przez wartość; wywołujący świadomie używa \texttt{std::move}. \\
\hline
\end{tabular}
\end{center}

\paragraph{Najważniejsze rozróżnienie.}
\begin{enumerate}
\item \texttt{std::move} wyraża zgodę na potraktowanie obiektu jako źródła przeniesienia, ale samo zasobów nie przenosi.
\item \texttt{std::forward<T>} zachowuje kategorię argumentu dedukowanego w szablonie.
\item Konstruktor przenoszący to zwykle operacja \texttt{T(T\&\&)}, ale może wykonywać dowolną operację zgodną z kontraktem typu.
\item \texttt{T\&\&} jako typ parametru nie zawsze jest referencją przekazującą.
\item Nazwany parametr typu \texttt{T\&\&} jest wyrażeniem \texttt{lvalue}.
\item Zwracanie przez wartość jest często najlepszym sposobem przekazywania wyniku.
\item Kiedy tylko się da, stosuj regułę zera zamiast własnych pięciu funkcji specjalnych.
\end{enumerate}

\paragraph{Pytania kontrolne.}
\begin{enumerate}
\item Dlaczego wywołanie \texttt{g(x)} w funkcji \texttt{f(T\&\& x)} traktuje \texttt{x} jako lvalue?
\item Kiedy \texttt{std::move(x)} kończy się kopiowaniem zamiast przeniesienia?
\item Czym różnią się \texttt{return T\{\};} i \texttt{return local;} w C++20?
\item Dlaczego przy powiększaniu \texttt{vector} konstruktor move z \texttt{noexcept} ma znaczenie?
\item Kiedy użyć \texttt{std::forward<T>(x)}, a kiedy \texttt{std::move(x)}?
\item Dlaczego metoda zwracająca \texttt{const std::string\&} nie powinna ujawniać referencji do składowej obiektu tymczasowego?
\end{enumerate}

\paragraph{Dalsza lektura.} \emph{cppreference}: \textit{Value categories},
\textit{Move constructors}, \textit{Reference declaration},
\textit{Copy elision}, \textit{std::move}, \textit{std::forward}.


\chapter[Semantyka przenoszenia w C++]{Semantyka przenoszenia, kategorie wartości i przekazywanie doskonałe}
\label{chap:move-semantics-cpp}

\section{Dlaczego wprowadzono semantykę przenoszenia?}

Rozważmy kontener przechowujący milion liczb. Skopiowanie go oznacza
zwykle przydzielenie nowej pamięci i skopiowanie każdego elementu.
Jeżeli dotychczasowy obiekt nie będzie już potrzebny, bardziej naturalna
jest operacja \emph{przekazania własności} jego zasobu zamiast kopiowania.
Od C++11 język udostępnia \emph{semantykę przenoszenia} (\textit{move
semantics}): programista może rozróżnić sytuację, w której obiekt źródłowy
ma pozostać niezmieniony, od sytuacji, w której jego zasoby wolno wykorzystać
ponownie.

\begin{lstlisting}
#include <string>
#include <utility>
#include <vector>

std::vector<std::string> a(100'000, "data");
auto b = a;                // copy: a remains intact

auto c = std::move(a);     // move may transfer a's allocation
// c now owns its elements; a is still valid, but its state
// after this move should not be assumed without a guarantee.
\end{lstlisting}

Nie oznacza to, że każde przeniesienie jest szybsze niż kopiowanie.
Dla typu \texttt{int} oba działania są w praktyce tożsame, a niektóre typy
kopiują część małych danych również podczas przenoszenia.
Kluczowa idea jest semantyczna: \emph{przeniesienie pozwala zmienić stan
źródła, a kopiowanie nie powinno go zmieniać}.

\paragraph{Zakres rozdziału.}
Zakładamy \textbf{C++20}. Kategorie wartości istnieją od początku języka,
referencje do r-wartości i \texttt{std::move} wprowadzono w C++11,
\emph{guaranteed copy elision} w C++17, a wiele przykładów poniżej
wykorzystuje koncepcje omówione w rozdziale~5.
Wszystkie przykłady zawierają komentarze angielskie, aby łatwo
przenosić je do projektów.

\section{Typ wyrażenia a kategoria wartości}
\label{sec:move-categories}

\textbf{Typ} i \textbf{kategoria wartości} są dwiema różnymi własnościami
wyrażenia. Typ mówi, jaki rodzaj danych otrzymujemy; kategoria mówi
m.in., czy wyrażenie identyfikuje istniejący obiekt i jak można
wiązać je z referencjami. Dokładnie trzy podstawowe kategorie to:
\emph{lvalue}, \emph{xvalue} i \emph{prvalue}.

\begin{center}
\begin{tabular}{|l|p{10.5cm}|}
\hline
\textbf{Kategoria} & \textbf{Intuicja i przykłady} \\
\hline
\texttt{lvalue} & Wyrażenie identyfikujące obiekt lub funkcję; np. \texttt{x}, \texttt{*ptr}, \texttt{++x} dla wbudowanego \texttt{int}. \\
\hline
\texttt{xvalue} & \emph{Expiring value}: wyrażenie identyfikujące obiekt, którego zasoby mogą być ponownie wykorzystane; np. \texttt{std::move(x)}. \\
\hline
\texttt{prvalue} & \emph{Pure rvalue}: wyrażenie obliczające wartość lub inicjalizujące obiekt; np. \texttt{42}, \texttt{x + 1}, \texttt{std::string("hi")}. \\
\hline
\end{tabular}
\end{center}

Pozostałe dwie nazwy oznaczają sumy kategorii:
\[
\underbrace{\mathrm{glvalue}}_{\mathrm{lvalue}\,\cup\,\mathrm{xvalue}},
\qquad
\underbrace{\mathrm{rvalue}}_{\mathrm{prvalue}\,\cup\,\mathrm{xvalue}}.
\]
Nie należy utożsamiać \texttt{rvalue} z tymczasowym obiektem ani
\texttt{lvalue} z wyrażeniem znajdującym się po lewej stronie przypisania.
Są to techniczne kategorie zdefiniowane przez język.

\subsection{Kilka wyrażeń o zaskakujących kategoriach}

\begin{lstlisting}
#include <string>
#include <utility>

int x = 7;
int* p = &x;

// x                 : lvalue
// *p                : lvalue
// (x)               : lvalue
// 7                 : prvalue
// x + 1             : prvalue
// std::move(x)      : xvalue
// std::string{"abc"}: prvalue

int& lref = x;
int&& rref = 9;

// lref              : lvalue
// rref              : lvalue (!) - it has a name
// std::move(rref)   : xvalue
\end{lstlisting}

To fundamentalna zasada: \textbf{nazwa zmiennej jest wyrażeniem lvalue,
nawet jeśli typ tej zmiennej to referencja rvalue \texttt{T\&\&}}.
Samo zadeklarowanie \texttt{T\&\&} nie oznacza, że każde użycie takiej
zmiennej będzie r-wartością.

\subsection{Badanie kategorii za pomocą \texttt{decltype((expr))}}

Dla prostego, nieujętego w nawiasy identyfikatora składnia
\texttt{decltype(x)} zwraca deklarowany typ. Dla ogólnego wyrażenia
\texttt{decltype((expr))} oddaje jego kategorię poprzez typ:
\[
\mathrm{decltype}((e))=
\left\{
\begin{array}{ll}
T\&, & \mbox{gdy } e \mbox{ jest lvalue typu } T,\\
T\&\&, & \mbox{gdy } e \mbox{ jest xvalue typu } T,\\
T, & \mbox{gdy } e \mbox{ jest prvalue typu } T.
\end{array}
\right.
\]

\begin{lstlisting}
#include <type_traits>
#include <utility>

int x = 10;
static_assert(std::is_same_v<decltype(x), int>);
static_assert(std::is_same_v<decltype((x)), int&>);
static_assert(std::is_same_v<decltype((std::move(x))), int&&>);
static_assert(std::is_same_v<decltype((x + 1)), int>);
\end{lstlisting}

\paragraph{Uwaga.} \texttt{decltype} stosuje szczególne reguły dla
nieujętych w nawiasy nazw i dostępu do składowych. Dlatego
\texttt{decltype(x)} oraz \texttt{decltype((x))} mogą zwrócić
różne typy. Nie jest to sprzeczność, lecz zamierzona cecha języka.

\subsection{Różnica między \texttt{auto} i \texttt{decltype(auto)}}

Zwykłe \texttt{auto} w deklaracji zmiennej najczęściej dedukuje typ
przez wartość i pomija referencję wynikającą z inicjalizatora.
Natomiast \texttt{decltype(auto)} stosuje reguły \texttt{decltype}
i może zachować kategorię wartości wyrażenia.

\begin{lstlisting}
#include <type_traits>

int n = 5;
auto copy = (n);           // int: independent value
decltype(auto) alias = (n); // int&: refers to n

static_assert(std::is_same_v<decltype(copy), int>);
static_assert(std::is_same_v<decltype(alias), int&>);

auto&& first = n;          // deduces int&
auto&& second = 5;         // deduces int&&
static_assert(std::is_same_v<decltype(first), int&>);
static_assert(std::is_same_v<decltype(second), int&&>);
\end{lstlisting}

W typie wyniku funkcji zapis \texttt{decltype(auto)} wymaga
ostrożności. Instrukcja \texttt{return (x);} może zwrócić
referencję do zmiennej lokalnej, która przestanie istnieć.
Zwykłe \texttt{auto} jest często bezpieczniejszym wyborem,
gdy funkcja ma zwrócić własny obiekt przez wartość.

\section{Referencje lvalue i rvalue}
\label{sec:move-references}

\begin{lstlisting}
#include <string>

std::string name = "Ada";
std::string& a = name;            // binds to mutable lvalue
const std::string& b = name;      // binds to lvalue
const std::string& c = "temporary"; // can bind to temporary
std::string&& d = std::string{"x"}; // binds to rvalue

// std::string& e = std::string{"x"}; // error
// std::string&& f = name;           // error
\end{lstlisting}

Referencja \texttt{T\&} umożliwia pracę z istniejącym, zwykle
nazwanym obiektem; \texttt{T\&\&} może wiązać się z r-wartością.
\texttt{const T\&} wiąże się zarówno z lvalue, jak i z rvalue,
nie dając prawa do modyfikowania obiektu poprzez tę referencję.

\subsection{Wybór przeciążenia}

\begin{lstlisting}
#include <iostream>
#include <string>
#include <utility>

void inspect(const std::string&) {
    std::cout << "const lvalue reference\n";
}
void inspect(std::string&&) {
    std::cout << "rvalue reference\n";
}

int main() {
    std::string text = "hello";
    inspect(text);             // const lvalue reference
    inspect(std::string{"x"}); // rvalue reference
    inspect(std::move(text));  // rvalue reference
}
\end{lstlisting}

Istnienie przeciążenia \texttt{f(T\&\&)} samo w sobie nie przesądza
o rzeczywistym przeniesieniu czegokolwiek. Funkcja może tylko
odczytać argument; \emph{move} jest operacją wykonywaną przez
konstruktor lub przypisanie, a nie przez mechanizm wiązania referencji.

\subsection{Wydłużenie życia obiektu tymczasowego}

W typowych deklaracjach lokalnych związanie obiektu tymczasowego
bezpośrednio z referencją może wydłużyć jego życie do końca życia
tej referencji:
\begin{lstlisting}
#include <string>

const std::string& a = std::string{"safe until a ends"};
std::string&& b = std::string{"safe until b ends"};
\end{lstlisting}
Nie wolno jednak uogólniać tej reguły na zwracanie referencji ani
na dowolne przekazywanie referencji dalej. Przykładowo:
\begin{lstlisting}
#include <string>

const std::string& broken() {
    return std::string{"dangling"}; // dangling reference
}
\end{lstlisting}
Zwrócona referencja nie utrzymuje tego obiektu przy życiu.
Jeżeli funkcja ma \emph{wytworzyć} wynik, zwykle powinna zwracać go
przez wartość.

\section{Co naprawdę robi \texttt{std::move}?}
\label{sec:move-std-move}

Nazwa jest nieco myląca: \texttt{std::move(x)} \textbf{nie przenosi
zasobów samodzielnie}. W przybliżeniu wykonuje rzutowanie do
referencji rvalue odpowiedniego typu. Dopiero dalszy konstruktor,
operator przypisania lub inna funkcja może potraktować wynik
jako obiekt nadający się do przeniesienia.

\begin{lstlisting}
#include <type_traits>
#include <utility>

template<class T>
constexpr std::remove_reference_t<T>&&
my_move(T&& value) noexcept {
    return static_cast<std::remove_reference_t<T>&&>(value);
}
\end{lstlisting}

Powyższa implementacja ilustruje istotę działania
\texttt{std::move}. W rzeczywistym kodzie korzystamy z funkcji
standardowej z nagłówka \texttt{<utility>}.

\begin{lstlisting}
#include <string>
#include <utility>

std::string first = "long text";
std::string second = std::move(first);
// A move constructor is selected for second.
// first is valid; do not assume its exact contents.

first = "reused";             // valid: assignment
first.clear();                // valid: clear()
\end{lstlisting}

\paragraph{Stan obiektu po przeniesieniu.}
Typy biblioteczne zwykle gwarantują stan \emph{poprawny, lecz
nieokreślony co do zawartości} (\textit{valid but unspecified}), o ile
kontrakt konkretnego typu nie obiecuje więcej. Takiego obiektu
można bezpiecznie zniszczyć, przypisać mu nową wartość albo
wykonać operację, której wymagania wstępne na to pozwalają.
Nie należy jednak \emph{zakładać}, że każdy przeniesiony
\texttt{std::string} jest pusty. Dla niektórych typów, np.
\texttt{std::unique\_ptr}, stan po przeniesieniu określono dokładniej.

\subsection{Problem obiektu \texttt{const}}

\begin{lstlisting}
#include <string>
#include <utility>

const std::string original = "cannot be emptied here";
std::string other = std::move(original);
// Typically: copy, not move (const string&& cannot bind to
// string's non-const move constructor).
\end{lstlisting}

\texttt{std::move(original)} ma w tym przypadku typ
\texttt{const std::string\&\&}. Typowy konstruktor przenoszący
przyjmuje \texttt{std::string\&\&}, więc nie może otrzymać
argumentu \texttt{const}. Kompilator wybiera konstruktor kopiujący
przyjmujący \texttt{const std::string\&}. \textbf{Samo wystąpienie
\texttt{std::move} w kodzie nie gwarantuje przeniesienia}.\footnote{Dla
typów niekopiowalnych ten sam zapis może skończyć się błędem
kompilacji.}

\section{Konstruktor przenoszący i przenoszący operator przypisania}
\label{sec:move-special-members}

W klasie \texttt{T} konstruktor przenoszący ma zwykle sygnaturę
\texttt{T(T\&\&)}, natomiast przenoszący operator przypisania:
\texttt{T\& operator=(T\&\&)}. Nie należy mylić ich z kopiowaniem:
\begin{lstlisting}
struct Item {
    Item() = default;
    Item(const Item&);            // copy constructor
    Item& operator=(const Item&); // copy assignment
    Item(Item&&) noexcept;        // move constructor
    Item& operator=(Item&&) noexcept; // move assignment
    ~Item();                      // destructor
};
\end{lstlisting}

\subsection{Przykład z obserwacją operacji}
\label{sec:move-trace}

\begin{lstlisting}
#include <iostream>
#include <utility>

struct Trace {
    Trace() { std::cout << "default\n"; }
    Trace(const Trace&) { std::cout << "copy ctor\n"; }
    Trace(Trace&&) noexcept { std::cout << "move ctor\n"; }
    Trace& operator=(const Trace&) {
        std::cout << "copy assign\n";
        return *this;
    }
    Trace& operator=(Trace&&) noexcept {
        std::cout << "move assign\n";
        return *this;
    }
};

int main() {
    Trace a;                    // default
    Trace b = a;                // copy ctor
    Trace c = std::move(a);     // move ctor
    b = c;                      // copy assign
    b = std::move(c);           // move assign
}
\end{lstlisting}

Konstruktor tworzy \emph{nowy} obiekt, a operator przypisania
zmienia już istniejący. W tym przykładzie \texttt{Trace} jedynie
wypisuje komunikaty. W klasie zarządzającej zasobem konstruktor
przenoszący powinien przekazać własność zasobu, a operator
przypisania musi dodatkowo poprawnie pozbyć się zasobu
posiadanego wcześniej przez cel.

\subsection{Pełny przykład: klasa samodzielnie zarządzająca tablicą}
\label{sec:move-raw-buffer}

To przykład \emph{edukacyjny}. W rzeczywistym projekcie zwykle
lepiej użyć standardowego kontenera, np. \texttt{std::vector<int>}.

\begin{lstlisting}
#include <algorithm>
#include <cstddef>
#include <utility>

class RawBuffer {
    std::size_t size_ = 0;
    int* data_ = nullptr;

public:
    explicit RawBuffer(std::size_t n)
        : size_(n), data_(n ? new int[n]{} : nullptr) {}

    ~RawBuffer() { delete[] data_; }

    RawBuffer(const RawBuffer& other)
        : RawBuffer(other.size_) {
        if (size_ != 0)
            std::copy_n(other.data_, size_, data_);
    }

    RawBuffer& operator=(const RawBuffer& other) {
        if (this != &other) {
            RawBuffer copy(other); // may throw; *this unchanged
            swap(copy);
        }
        return *this;
    }

    RawBuffer(RawBuffer&& other) noexcept
        : size_(std::exchange(other.size_, 0)),
          data_(std::exchange(other.data_, nullptr)) {}

    RawBuffer& operator=(RawBuffer&& other) noexcept {
        if (this != &other) {
            delete[] data_;
            size_ = std::exchange(other.size_, 0);
            data_ = std::exchange(other.data_, nullptr);
        }
        return *this;
    }

    void swap(RawBuffer& other) noexcept {
        using std::swap;
        swap(size_, other.size_);
        swap(data_, other.data_);
    }

    std::size_t size() const noexcept { return size_; }
    int& operator[](std::size_t i) { return data_[i]; }
    const int& operator[](std::size_t i) const { return data_[i]; }
};
\end{lstlisting}

\texttt{std::exchange(obj, value)} zwraca starą wartość \texttt{obj}
i zastępuje ją wartością \texttt{value}. Konstruktor przenoszący
przejmuje więc wskaźnik i rozmiar, a źródło ustawia w jawnie
określony stan pusty. Operator przenoszącego przypisania
usuwa dotychczasową tablicę celu, przejmuje nową i zabezpiecza
przypadek \texttt{b = std::move(b)}. W implementacji kopiowania
widać z kolei \emph{silną gwarancję wyjątkową}: jeśli utworzenie
kopii nie powiedzie się, pierwotny obiekt pozostaje bez zmian.

\section{Reguła pięciu, reguła trzech i reguła zera}
\label{sec:move-rule-of-five}

Klasy zarządzające zasobem często muszą kontrolować pięć operacji:
(1) destruktor, (2) konstruktor kopiujący, (3) kopiujący operator
przypisania, (4) konstruktor przenoszący oraz (5) przenoszący
operator przypisania. To \emph{reguła pięciu}.

Historyczna \emph{reguła trzech} odnosiła się do destruktora,
konstruktora kopiującego i kopiującego operatora przypisania,
zanim język udostępnił move semantics.

\paragraph{Ważne: deklarowanie destruktora może blokować generowanie
konstruktora przenoszącego.}

\begin{lstlisting}
#include <memory>
#include <type_traits>

struct WithDestructor {
    std::unique_ptr<int> p;
    ~WithDestructor() = default;
};

static_assert(!std::is_move_constructible_v<WithDestructor>);

struct WithMove {
    std::unique_ptr<int> p;
    ~WithMove() = default;
    WithMove() = default;
    WithMove(WithMove&&) noexcept = default;
    WithMove& operator=(WithMove&&) noexcept = default;
};
static_assert(std::is_move_constructible_v<WithMove>);
\end{lstlisting}

W \texttt{WithDestructor} użytkownik zadeklarował destruktor,
więc kompilator nie deklaruje automatycznie operacji przenoszących.
Kopiowanie jest niedostępne z powodu \texttt{unique\_ptr}.
Samo \texttt{= default} przy destruktorze nie usuwa tej reguły.

\subsection{Najlepsza praktyka: reguła zera}

\emph{Reguła zera} mówi, że klasy, które nie są bezpośrednio
odpowiedzialne za posiadanie surowych zasobów, zwykle nie powinny
same definiować żadnej z pięciu operacji.

\begin{lstlisting}
#include <string>
#include <vector>

struct Dataset {
    std::string title;
    std::vector<double> values;
};

Dataset original{"results", {1.0, 2.0, 3.0}};
Dataset copy = original;           // deep copy of vector data
Dataset moved = std::move(original); // member-wise move
\end{lstlisting}

To kontenery standardowe odpowiadają za przydział i zwalnianie
pamięci. Automatycznie generowane operacje \texttt{Dataset}
wykonują odpowiednie operacje na jego składowych. \textbf{Zalecenie:}
zacznij od reguły zera; własny konstruktor przenoszący pisz
jedynie wtedy, gdy zachowanie domyślne nie jest właściwe.

\section{Przenoszenie a wyjątki: \texttt{noexcept}}
\label{sec:move-noexcept}

Kontenery takie jak \texttt{std::vector<T>} przy realokacji muszą
przenieść lub skopiować dotychczasowe elementy do nowego obszaru.
Jeśli przenoszenie \texttt{T} może rzucać wyjątki, a \texttt{T}
można kopiować, biblioteka może wybrać kopiowanie, by zachować
silniejsze gwarancje wyjątkowe. Właśnie dlatego poprawne operacje
przenoszące, które nie mogą rzucać wyjątków, warto deklarować
\texttt{noexcept}.

\begin{lstlisting}
#include <type_traits>
#include <utility>

struct Resource {
    Resource() = default;
    Resource(const Resource&) = default;
    Resource(Resource&&) noexcept = default;
};

static_assert(std::is_nothrow_move_constructible_v<Resource>);

// move_if_noexcept can provide an lvalue reference if
// moving might throw and a copy is available.
\end{lstlisting}

\texttt{std::move\_if\_noexcept(x)} zwraca odpowiednio
\texttt{T\&\&} albo \texttt{const T\&}, zależnie od własności
konstruktorów danego typu. To narzędzie używane przy implementowaniu
kontenerów, a nie uniwersalny zamiennik \texttt{std::move}.

\paragraph{Ostrożnie z \texttt{noexcept}.}
Wyjście wyjątku poza funkcję zadeklarowaną jako
\texttt{noexcept} wywołuje \texttt{std::terminate}.
Dlatego bezwyjątkowość musi wynikać z rzeczywistych operacji składowych.

\section{Zwijanie referencji (\textit{reference collapsing})}
\label{sec:move-ref-collapse}

Gdy referencje występują pośrednio poprzez aliasy lub
podstawienie argumentów szablonu, stosuje się \emph{zwijanie
referencji}. Wszystkie przypadki przedstawia tabela:

\begin{center}
\begin{tabular}{|c|c|c|}
\hline
\textbf{Typ pośredni} & \textbf{Zapis} & \textbf{Wynik} \\
\hline
\texttt{T\&} i \texttt{\&} & \texttt{T\& \&} & \texttt{T\&} \\
\texttt{T\&} i \texttt{\&\&} & \texttt{T\& \&\&} & \texttt{T\&} \\
\texttt{T\&\&} i \texttt{\&} & \texttt{T\&\& \&} & \texttt{T\&} \\
\texttt{T\&\&} i \texttt{\&\&} & \texttt{T\&\& \&\&} & \texttt{T\&\&} \\
\hline
\end{tabular}
\end{center}

Krótko: jeżeli \emph{choć jeden} składnik jest \texttt{\&},
wynik jest \texttt{\&}; tylko \texttt{\&\&} połączone
z \texttt{\&\&} daje \texttt{\&\&}.

\begin{lstlisting}
#include <type_traits>

using L = int&;
using R = int&&;

static_assert(std::is_same_v<L&, int&>);
static_assert(std::is_same_v<L&&, int&>);
static_assert(std::is_same_v<R&, int&>);
static_assert(std::is_same_v<R&&, int&&>);
\end{lstlisting}

Zwijanie referencji jest podstawą \emph{perfect forwarding}:
dzięki niemu jeden szablon może przyjmować zarówno lvalue,
jak i rvalue i przekazywać tę różnicę dalej.

\section{Referencje przekazujące i dedukcja \texttt{T\&\&}}
\label{sec:move-forwarding-ref}

Zapis \texttt{T\&\&} ma dwa znaczenia w zależności od kontekstu.
Jeśli \texttt{T} jest dedukowanym, niekwalifikowanym parametrem
typu \emph{tego samego} szablonu funkcji, \texttt{T\&\&}
to \emph{referencja przekazująca} (\textit{forwarding reference}).
W pozostałych typowych kontekstach jest to po prostu
referencja rvalue.

\begin{lstlisting}
#include <string>

template<class T>
void forwarding(T&& x);         // forwarding reference

template<class T>
void not_forwarding(const T&& x); // rvalue reference

void fixed(std::string&& x);    // rvalue reference

template<class T>
struct Holder {
    void receive(T&& x);        // NOT forwarding: T comes
                               // from the class template

    template<class U>
    void relay(U&& x);          // forwarding: U is deduced
};
\end{lstlisting}

Dla wywołania \texttt{forwarding(x)} z lvalue typu \texttt{int}
kompilator dedukuje \texttt{T = int\&}, zatem parametr ma po
zwinięciu typ \texttt{int\&}. Dla \texttt{forwarding(5)}
dedukuje \texttt{T = int}, a parametr ma typ \texttt{int\&\&}.

\begin{lstlisting}
#include <type_traits>

template<class T>
void check(T&& x) {
    if constexpr (std::is_lvalue_reference_v<T>) {
        // Called with an lvalue.
    } else {
        // Called with an rvalue.
    }
}
\end{lstlisting}

\paragraph{Drobne wyjątki.} \texttt{auto\&\&} zwykle również
pozwala na przekazywanie kategorii, lecz dedukcja \texttt{auto}
ma szczególne reguły dla list inicjalizacyjnych.
\texttt{const T\&\&} nie jest referencją przekazującą.

\section{Perfect forwarding i \texttt{std::forward}}
\label{sec:move-perfect-forwarding}

Przez \emph{perfect forwarding} rozumiemy przekazanie argumentu
do następnej funkcji przy zachowaniu kategorii wartości,
a także odpowiednich kwalifikatorów typu.

\begin{lstlisting}
#include <iostream>
#include <utility>

void target(int&)  { std::cout << "lvalue\n"; }
void target(int&&) { std::cout << "rvalue\n"; }

template<class T>
void incorrect(T&& value) {
    target(value); // value is a named variable: always lvalue
}

template<class T>
void correct(T&& value) {
    target(std::forward<T>(value));
}

int main() {
    int x = 1;
    incorrect(2); // lvalue (!)
    correct(x);   // lvalue
    correct(2);   // rvalue
}
\end{lstlisting}

W przybliżeniu \texttt{std::forward<T>(value)} zwraca
lvalue, jeżeli \texttt{T} wydedukowano jako referencję lvalue,
i xvalue w przeciwnym razie. \texttt{std::move(value)} zawsze
wymusza kategorię xvalue: w funkcji przekazującej byłoby to
błędem, ponieważ niepotrzebnie pozwalałoby na przeniesienie
zasobu z argumentu przekazanego jako lvalue.

\subsection{Przykład: fabryka obiektów}

\begin{lstlisting}
#include <memory>
#include <string>
#include <utility>

struct Person {
    std::string name;
    int age;

    Person(std::string n, int a)
        : name(std::move(n)), age(a) {}
};

template<class T, class... Args>
std::unique_ptr<T> make_object(Args&&... args) {
    return std::unique_ptr<T>(
        new T(std::forward<Args>(args)...)
    );
}

int main() {
    std::string first = "Ada";
    auto p = make_object<Person>(first, 37); // first copied
    auto q = make_object<Person>(std::move(first), 38);
}
\end{lstlisting}

To edukacyjna wersja fabryki. W produkcyjnym kodzie
\texttt{std::make\_unique<T>(...)} robi to samo w prostszy
sposób. W drugim wywołaniu użytkownik świadomie pozwolił na
przeniesienie \texttt{first}; w pierwszym pozostaje on
nienaruszony. \texttt{std::forward<Args>(args)...} rozwija pakiet
argumentów omówiony w rozdziale~4.

\subsection{Nie przekazuj tego samego zasobu dwa razy bez zastanowienia}

\begin{lstlisting}
#include <utility>

template<class T>
void bad_double_forward(T&& x) {
    consume(std::forward<T>(x));
    consume(std::forward<T>(x)); // may use moved-from object
}
\end{lstlisting}

Ten fragment jest \emph{schematem błędu}: zakładamy, że istnieje
funkcja \texttt{consume}. Jeśli pierwsze wywołanie wykorzysta
argument rvalue przez przeniesienie, drugie otrzyma już obiekt
po przeniesieniu. Referencję przekazującą traktuj zatem jako
przeznaczoną zwykle do \emph{jednego} przekazania dalej.

\section{Perfect forwarding a przeciążenia i szablony klas}

Konstruktor z referencją przekazującą potrafi być bardzo zachłanny
podczas rozstrzygania przeciążeń. Może np. wygrać z konstruktorem
kopiującym przy argumentach nie-\texttt{const}.

\begin{lstlisting}
#include <string>
#include <utility>

struct Named {
    std::string name;

    Named(const Named&) = default;

    template<class U>
    explicit Named(U&& u)
        : name(std::forward<U>(u)) {}
};

// Named a{"text"};
// Named b(a);  // may select the constructor template,
//              // then fail to initialize std::string from Named&
\end{lstlisting}

Zwykle lepiej zacząć od zwykłych przeciążeń albo konstruktora
przyjmującego konkretny typ. Jeżeli szablon jest potrzebny,
można ograniczyć go koncepcją:

\begin{lstlisting}
#include <concepts>
#include <string>
#include <type_traits>
#include <utility>

struct SafeNamed {
    std::string name;

    SafeNamed(const SafeNamed&) = default;
    SafeNamed(SafeNamed&&) = default;

    template<class U>
      requires (!std::same_as<std::remove_cvref_t<U>, SafeNamed>
                && std::constructible_from<std::string, U&&>)
    explicit SafeNamed(U&& u)
        : name(std::forward<U>(u)) {}
};
\end{lstlisting}

Takie powiązanie rozdziału o perfect forwarding z koncepcjami
z rozdziału~5 jest bardzo przydatne
w bibliotekach generycznych.

\section{Copy elision: kiedy nie ma ani kopii, ani przeniesienia?}
\label{sec:move-elision}

Prosty model \emph{zwrócenie przez wartość = kopia lub move}
jest niepoprawny w nowoczesnym C++. Od C++17 pewne wyrażenia
\texttt{prvalue} inicjalizują wynik \emph{bezpośrednio w miejscu
docelowym}. Nie istnieje dodatkowy obiekt, który trzeba by
skopiować lub przenieść.

\begin{lstlisting}
struct NonMovable {
    NonMovable() = default;
    NonMovable(const NonMovable&) = delete;
    NonMovable(NonMovable&&) = delete;
};

NonMovable create() {
    return NonMovable{}; // guaranteed elision, C++17+
}

int main() {
    NonMovable x = create(); // OK in C++17+
}
\end{lstlisting}

Nie jest to zwykła optymalizacja, którą kompilator \emph{może}
wyłączyć: w pokazanym przypadku wynika to z semantyki języka.

\subsection{NRVO a gwarantowana elizja}

\begin{lstlisting}
#include <string>

std::string first() {
    return std::string{"hello"}; // mandatory direct construction
}

std::string second() {
    std::string s = "hello";
    return s;                    // NRVO may occur
}
\end{lstlisting}

\textit{Named Return Value Optimization} (NRVO) dotyczy
zwracania nazwanej lokalnej zmiennej. W C++20 kompilator \emph{może}
pominąć utworzenie dodatkowego obiektu. Jeśli NRVO nie nastąpi,
wyrażenie \texttt{s} w instrukcji \texttt{return} może
uczestniczyć w niejawnej próbie przeniesienia.

\paragraph{Nie pisz automatycznie \texttt{return std::move(s);}.}

\begin{lstlisting}
std::string make_text() {
    std::string s = "result";
    return s;             // allows NRVO
}

std::string less_good() {
    std::string s = "result";
    return std::move(s);  // inhibits NRVO; move may be required
}
\end{lstlisting}

\texttt{std::move(s)} zmienia wynik w xvalue i tym samym
nie spełnia warunków NRVO dotyczących nazwy obiektu.
Jeśli funkcja zwraca własną lokalną zmienną przez wartość,
zazwyczaj najlepsze jest \texttt{return s;}.

\section{Argumenty funkcji: kiedy kopiować, kiedy przenosić?}
\label{sec:move-api}

Wybór sposobu przyjmowania argumentu jest częścią kontraktu API.
Nie chodzi jedynie o wydajność, ale przede wszystkim o własność
i intencję.

\begin{center}
\begin{tabular}{|p{3.4cm}|p{4.3cm}|p{4.7cm}|}
\hline
\textbf{Parametr} & \textbf{Znaczenie} & \textbf{Kiedy używać} \\
\hline
\texttt{const T\&} & Odczyt bez przejmowania własności. & Duże obiekty, gdy wystarczy tylko odczyt. \\
\hline
\texttt{T\&} & Modyfikacja obiektu należącego do wywołującego. & Parametry wyjściowe i operacje in-place. \\
\hline
\texttt{T} & Własna wartość; może być skopiowana albo przeniesiona. & Małe typy i funkcje mające zachować wartość. \\
\hline
\texttt{T\&\&} & Dostęp do argumentu będącego r-wartością. & Przeciążenia zależne od kategorii oraz operacje konsumujące. \\
\hline
\texttt{T\&\&} z dedukcją \texttt{T} & Referencja przekazująca. & Opakowania i fabryki zachowujące kategorię argumentów. \\
\hline
\end{tabular}
\end{center}

\subsection{Idiom \emph{pass by value, then move}}

\begin{lstlisting}
#include <string>
#include <utility>

class User {
    std::string name_;
public:
    explicit User(std::string name)
        : name_(std::move(name)) {}
};

std::string a = "Alice";
User first(a);                    // copy into parameter, then move
User second(std::string{"Bob"});  // parameter from prvalue, then move
\end{lstlisting}

To dobre rozwiązanie, gdy funkcja \emph{zawsze} potrzebuje własnej
kopii/przeniesienia argumentu. Nie jest jednak magicznie
najszybsze w każdym przypadku: dla lvalue może wymagać kopiowania
i jeszcze jednego przeniesienia, a duże obiekty niemożliwe do
przeniesienia mogą wymagać osobnego projektu API.

\subsection{Przekazywanie własności \texttt{unique\_ptr}}

\begin{lstlisting}
#include <memory>
#include <utility>

void adopt(std::unique_ptr<int> p) {
    // Ownership transferred into this function.
}

int main() {
    auto p = std::make_unique<int>(42);
    adopt(std::move(p));
    // p == nullptr after moving this unique_ptr
}
\end{lstlisting}

Przekazanie \texttt{unique\_ptr} przez wartość oznacza przejęcie
własności. Jeżeli funkcja ma jedynie korzystać z obiektu,
zwykle lepiej przyjąć \texttt{T\&} lub \texttt{const T\&}.
Temat własności zasobów omówiono w rozdziale~3.

\section{Rvalue-qualified methods i świadome udostępnianie zasobów}
\label{sec:move-ref-qualified}

C++ pozwala ograniczać metodę do lvalue albo rvalue odbiorcy,
dopisując \texttt{\&} lub \texttt{\&\&} po liście parametrów.
Można to wykorzystać np. do udostępnienia zawartości obiektu
bez kopiowania, gdy całe opakowanie jest tymczasowe.

\begin{lstlisting}
#include <string>
#include <utility>

class Message {
    std::string data_;
public:
    explicit Message(std::string s) : data_(std::move(s)) {}

    const std::string& text() const & noexcept {
        return data_;              // lvalue owner: borrow
    }
    std::string text() && noexcept {
        return std::move(data_);   // rvalue owner: consume
    }
};

Message m{"hello"};
const std::string& borrowed = m.text();
std::string owned = Message{"world"}.text();
\end{lstlisting}

Wariant \texttt{\&\&} może bezpiecznie oddać wartość przez
przeniesienie. Gdybyśmy zwracali w nim referencję do składowej
tymczasowego obiektu, powstałaby referencja wisząca.
Projektując takie API, warto rozważyć również wywołania na
\texttt{const} r-wartościach i ich wymagane zachowanie.

\section{Kontenery, iteratory przenoszące i wzorce użytkowe}

\subsection{Dodawanie elementów do \texttt{std::vector}}

\begin{lstlisting}
#include <string>
#include <utility>
#include <vector>

std::vector<std::string> words;
std::string word = "one";

words.push_back(word);             // copy from lvalue
words.push_back(std::move(word));  // move from lvalue cast to xvalue
words.emplace_back(3, 'x');       // construct "xxx" in place
\end{lstlisting}

Metoda \texttt{emplace\_back} konstruuje nowy element z podanych
argumentów. Nie musi być szybsza od \texttt{push\_back}.
Gotowy obiekt można przekazać z pomocą \texttt{std::move}.
Podczas realokacji kontener nadal może kopiować lub przenosić
wcześniej dodane elementy.

\subsection{Przenoszenie całego zakresu}

\begin{lstlisting}
#include <algorithm>
#include <iterator>
#include <string>
#include <vector>

std::vector<std::string> source{"a", "b"};
std::vector<std::string> dest;

std::move(source.begin(), source.end(),
          std::back_inserter(dest)); // algorithm std::move

// source still contains two valid strings, but their
// contents are not assumed.
\end{lstlisting}

Nazwa \texttt{std::move} występuje również dla algorytmu z nagłówka
\texttt{<algorithm>}. W przeciwieństwie do pojedynczego
\texttt{std::move(x)} z nagłówka \texttt{<utility>}, algorytm
przypisuje kolejno elementy z zakresu docelowego, korzystając
z możliwości ich przenoszenia.

\section{Typowe błędy i pułapki}
\label{sec:move-errors}

\subsection{Błąd 1: przekonanie, że \texttt{std::move} niszczy obiekt}

\texttt{std::move} jest rzutowaniem, a nie destruktorem. Obiekt
źródłowy nadal istnieje aż do końca własnego czasu życia.
Może zmienić stan, jeśli na jego podstawie wykonano
rzeczywistą operację przeniesienia.

\subsection{Błąd 2: \texttt{std::move} z argumentu przekazanego jako \texttt{const T\&}}

\begin{lstlisting}
#include <string>
#include <utility>

std::string ineffective(const std::string& x) {
    return std::move(x); // copies for std::string
}
\end{lstlisting}

Jeżeli funkcja potrzebuje zwrócić niezależną kopię, czytelniej
napisać \texttt{return x;}. Jeśli ma konsumować argument,
powinna otrzymać wartość lub odpowiednią modyfikowalną
referencję. Nie wolno usuwać \texttt{const} przez rzutowanie,
by wymusić przeniesienie z oryginalnie stałego obiektu.

\subsection{Błąd 3: użycie nazwanego \texttt{T\&\&} bez \texttt{std::move}}

\begin{lstlisting}
#include <memory>
#include <utility>

void take(std::unique_ptr<int>&& p) {
    // auto q = p;            // error: p is an lvalue
    auto q = std::move(p);    // explicit ownership transfer
}
\end{lstlisting}

\subsection{Błąd 4: przechowywanie referencji do argumentu tymczasowego}

\begin{lstlisting}
#include <string>
#include <string_view>

std::string_view dangerous() {
    return std::string{"temporary"}; // dangling view
}
\end{lstlisting}

\texttt{std::string\_view} jest widokiem, a nie właścicielem
pamięci. Semantyka przenoszenia nie rozwiązuje problemów
związanych z czasem życia. Gdy wynik funkcji ma być własnością
wywołującego, zwracaj \texttt{std::string} przez wartość.

\subsection{Błąd 5: przenoszenie nie gwarantuje zachowania wskaźników}

Przeniesienie kontenera może spowodować unieważnienie wcześniej
zachowanych iteratorów, referencji lub wskaźników zależnie od
typu kontenera, operacji i alokatora. Nie zakładaj ogólnej
zasady, że po \texttt{std::move} adresy danych zawsze pozostają
takie same. Sprawdź dokumentację odpowiedniego typu.

\subsection{Błąd 6: nadmierne stosowanie \texttt{std::forward}}

\texttt{std::forward<T>} ma sens głównie tam, gdzie \texttt{T}
pochodzi z dedukcji dla referencji przekazującej. Przy zwykłej
lokalnej zmiennej przeznaczonej do przekazania własności
należy używać \texttt{std::move}; przy zwracaniu lokalnej
zmiennej przez wartość zazwyczaj żadnego z nich.

\section{Przykład przekrojowy: kolejka zadań}
\label{sec:move-project}

Rozważmy prosty bufor zadań przechowujący funkcje do
późniejszego wykonania. Użyjemy \texttt{std::unique\_ptr},
aby zadanie jednoznacznie posiadało zasób. Poniższy przykład
jest niezależny od wielowątkowości; \emph{kolejka zadań}
nie uruchamia niczego w osobnym wątku.

\begin{lstlisting}
#include <functional>
#include <iostream>
#include <memory>
#include <utility>
#include <vector>

class Tasks {
    std::vector<std::function<void()>> items_;
public:
    template<class F>
    void add(F&& f) {
        items_.emplace_back(std::forward<F>(f));
    }

    void run() {
        for (auto& f : items_) f();
        items_.clear();
    }
};

int main() {
    Tasks tasks;
    auto value = std::make_shared<int>(10);
    tasks.add([value] { std::cout << *value << '\n'; });
    tasks.run();
}
\end{lstlisting}

\texttt{std::function<void()>} w C++20 wymaga, aby przechowywany
obiekt wywoływalny był kopiowalny. Dlatego w przykładzie lambda
przechowuje \texttt{shared\_ptr}. Gdybyśmy zamiast niego
przechwycili przez przeniesienie \texttt{unique\_ptr}, lambda
stałaby się niekopiowalna i taki obiekt nie spełniłby
wymagań \texttt{std::function} w C++20. C++23 wprowadza
\texttt{std::move\_only\_function}, która rozwiązuje ten problem.
To pokazuje, że prawidłowy wybór sposobu przenoszenia zależy
od kontraktów pozostałych typów.

\section{Praktyczna ściąga i zalecenia projektowe}

\begin{center}
\begin{tabular}{|p{4.0cm}|p{8.4cm}|}
\hline
\textbf{Sytuacja} & \textbf{Zalecenie} \\
\hline
Zwracam utworzony obiekt & Zwróć go przez wartość; korzystaj z elizji lub NRVO. \\
\hline
Przenoszę zasób z lokalnego obiektu & Użyj \texttt{std::move(x)} w miejscu przekazania własności. \\
\hline
Piszę wrapper o argumentach \texttt{Args\&\&...} & Wykorzystaj \texttt{std::forward<Args>(args)...}. \\
\hline
Klasa ma \texttt{vector}/\texttt{string} jako składowe & Preferuj regułę zera i generowane operacje. \\
\hline
Klasa posiada goły wskaźnik do alokacji & Preferuj zmianę projektu na RAII; jeśli zarządzasz nim ręcznie, rozważ wszystkie operacje specjalne. \\
\hline
Kontener przenosi elementy przy realokacji & Właściwy, bezpieczny \texttt{noexcept} w konstruktorze move pomaga w gwarancjach wyjątkowych. \\
\hline
Argument będzie tylko odczytywany & Zwykle \texttt{const T\&} dla dużego \texttt{T}, a małe typy przez wartość. \\
\hline
Funkcja przejmuje \texttt{unique\_ptr} & Przyjmij go przez wartość; wywołujący świadomie używa \texttt{std::move}. \\
\hline
\end{tabular}
\end{center}

\paragraph{Najważniejsze rozróżnienie.}
\begin{enumerate}
\item \texttt{std::move} wyraża zgodę na potraktowanie obiektu jako źródła przeniesienia, ale samo zasobów nie przenosi.
\item \texttt{std::forward<T>} zachowuje kategorię argumentu dedukowanego w szablonie.
\item Konstruktor przenoszący to zwykle operacja \texttt{T(T\&\&)}, ale może wykonywać dowolną operację zgodną z kontraktem typu.
\item \texttt{T\&\&} jako typ parametru nie zawsze jest referencją przekazującą.
\item Nazwany parametr typu \texttt{T\&\&} jest wyrażeniem \texttt{lvalue}.
\item Zwracanie przez wartość jest często najlepszym sposobem przekazywania wyniku.
\item Kiedy tylko się da, stosuj regułę zera zamiast własnych pięciu funkcji specjalnych.
\end{enumerate}

\paragraph{Pytania kontrolne.}
\begin{enumerate}
\item Dlaczego wywołanie \texttt{g(x)} w funkcji \texttt{f(T\&\& x)} traktuje \texttt{x} jako lvalue?
\item Kiedy \texttt{std::move(x)} kończy się kopiowaniem zamiast przeniesienia?
\item Czym różnią się \texttt{return T\{\};} i \texttt{return local;} w C++20?
\item Dlaczego przy powiększaniu \texttt{vector} konstruktor move z \texttt{noexcept} ma znaczenie?
\item Kiedy użyć \texttt{std::forward<T>(x)}, a kiedy \texttt{std::move(x)}?
\item Dlaczego metoda zwracająca \texttt{const std::string\&} nie powinna ujawniać referencji do składowej obiektu tymczasowego?
\end{enumerate}

\paragraph{Dalsza lektura.} \emph{cppreference}: \textit{Value categories},
\textit{Move constructors}, \textit{Reference declaration},
\textit{Copy elision}, \textit{std::move}, \textit{std::forward}.


\chapter[Semantyka przenoszenia w C++]{Semantyka przenoszenia, kategorie wartości i przekazywanie doskonałe}
\label{chap:move-semantics-cpp}

\section{Dlaczego wprowadzono semantykę przenoszenia?}

Rozważmy kontener przechowujący milion liczb. Skopiowanie go oznacza
zwykle przydzielenie nowej pamięci i skopiowanie każdego elementu.
Jeżeli dotychczasowy obiekt nie będzie już potrzebny, bardziej naturalna
jest operacja \emph{przekazania własności} jego zasobu zamiast kopiowania.
Od C++11 język udostępnia \emph{semantykę przenoszenia} (\textit{move
semantics}): programista może rozróżnić sytuację, w której obiekt źródłowy
ma pozostać niezmieniony, od sytuacji, w której jego zasoby wolno wykorzystać
ponownie.

\begin{lstlisting}
#include <string>
#include <utility>
#include <vector>

std::vector<std::string> a(100'000, "data");
auto b = a;                // copy: a remains intact

auto c = std::move(a);     // move may transfer a's allocation
// c now owns its elements; a is still valid, but its state
// after this move should not be assumed without a guarantee.
\end{lstlisting}

Nie oznacza to, że każde przeniesienie jest szybsze niż kopiowanie.
Dla typu \texttt{int} oba działania są w praktyce tożsame, a niektóre typy
kopiują część małych danych również podczas przenoszenia.
Kluczowa idea jest semantyczna: \emph{przeniesienie pozwala zmienić stan
źródła, a kopiowanie nie powinno go zmieniać}.

\paragraph{Zakres rozdziału.}
Zakładamy \textbf{C++20}. Kategorie wartości istnieją od początku języka,
referencje do r-wartości i \texttt{std::move} wprowadzono w C++11,
\emph{guaranteed copy elision} w C++17, a wiele przykładów poniżej
wykorzystuje koncepcje omówione w rozdziale~5.
Wszystkie przykłady zawierają komentarze angielskie, aby łatwo
przenosić je do projektów.

\section{Typ wyrażenia a kategoria wartości}
\label{sec:move-categories}

\textbf{Typ} i \textbf{kategoria wartości} są dwiema różnymi własnościami
wyrażenia. Typ mówi, jaki rodzaj danych otrzymujemy; kategoria mówi
m.in., czy wyrażenie identyfikuje istniejący obiekt i jak można
wiązać je z referencjami. Dokładnie trzy podstawowe kategorie to:
\emph{lvalue}, \emph{xvalue} i \emph{prvalue}.

\begin{center}
\begin{tabular}{|l|p{10.5cm}|}
\hline
\textbf{Kategoria} & \textbf{Intuicja i przykłady} \\
\hline
\texttt{lvalue} & Wyrażenie identyfikujące obiekt lub funkcję; np. \texttt{x}, \texttt{*ptr}, \texttt{++x} dla wbudowanego \texttt{int}. \\
\hline
\texttt{xvalue} & \emph{Expiring value}: wyrażenie identyfikujące obiekt, którego zasoby mogą być ponownie wykorzystane; np. \texttt{std::move(x)}. \\
\hline
\texttt{prvalue} & \emph{Pure rvalue}: wyrażenie obliczające wartość lub inicjalizujące obiekt; np. \texttt{42}, \texttt{x + 1}, \texttt{std::string("hi")}. \\
\hline
\end{tabular}
\end{center}

Pozostałe dwie nazwy oznaczają sumy kategorii:
\[
\underbrace{\mathrm{glvalue}}_{\mathrm{lvalue}\,\cup\,\mathrm{xvalue}},
\qquad
\underbrace{\mathrm{rvalue}}_{\mathrm{prvalue}\,\cup\,\mathrm{xvalue}}.
\]
Nie należy utożsamiać \texttt{rvalue} z tymczasowym obiektem ani
\texttt{lvalue} z wyrażeniem znajdującym się po lewej stronie przypisania.
Są to techniczne kategorie zdefiniowane przez język.

\subsection{Kilka wyrażeń o zaskakujących kategoriach}

\begin{lstlisting}
#include <string>
#include <utility>

int x = 7;
int* p = &x;

// x                 : lvalue
// *p                : lvalue
// (x)               : lvalue
// 7                 : prvalue
// x + 1             : prvalue
// std::move(x)      : xvalue
// std::string{"abc"}: prvalue

int& lref = x;
int&& rref = 9;

// lref              : lvalue
// rref              : lvalue (!) - it has a name
// std::move(rref)   : xvalue
\end{lstlisting}

To fundamentalna zasada: \textbf{nazwa zmiennej jest wyrażeniem lvalue,
nawet jeśli typ tej zmiennej to referencja rvalue \texttt{T\&\&}}.
Samo zadeklarowanie \texttt{T\&\&} nie oznacza, że każde użycie takiej
zmiennej będzie r-wartością.

\subsection{Badanie kategorii za pomocą \texttt{decltype((expr))}}

Dla prostego, nieujętego w nawiasy identyfikatora składnia
\texttt{decltype(x)} zwraca deklarowany typ. Dla ogólnego wyrażenia
\texttt{decltype((expr))} oddaje jego kategorię poprzez typ:
\[
\mathrm{decltype}((e))=
\left\{
\begin{array}{ll}
T\&, & \mbox{gdy } e \mbox{ jest lvalue typu } T,\\
T\&\&, & \mbox{gdy } e \mbox{ jest xvalue typu } T,\\
T, & \mbox{gdy } e \mbox{ jest prvalue typu } T.
\end{array}
\right.
\]

\begin{lstlisting}
#include <type_traits>
#include <utility>

int x = 10;
static_assert(std::is_same_v<decltype(x), int>);
static_assert(std::is_same_v<decltype((x)), int&>);
static_assert(std::is_same_v<decltype((std::move(x))), int&&>);
static_assert(std::is_same_v<decltype((x + 1)), int>);
\end{lstlisting}

\paragraph{Uwaga.} \texttt{decltype} stosuje szczególne reguły dla
nieujętych w nawiasy nazw i dostępu do składowych. Dlatego
\texttt{decltype(x)} oraz \texttt{decltype((x))} mogą zwrócić
różne typy. Nie jest to sprzeczność, lecz zamierzona cecha języka.

\subsection{Różnica między \texttt{auto} i \texttt{decltype(auto)}}

Zwykłe \texttt{auto} w deklaracji zmiennej najczęściej dedukuje typ
przez wartość i pomija referencję wynikającą z inicjalizatora.
Natomiast \texttt{decltype(auto)} stosuje reguły \texttt{decltype}
i może zachować kategorię wartości wyrażenia.

\begin{lstlisting}
#include <type_traits>

int n = 5;
auto copy = (n);           // int: independent value
decltype(auto) alias = (n); // int&: refers to n

static_assert(std::is_same_v<decltype(copy), int>);
static_assert(std::is_same_v<decltype(alias), int&>);

auto&& first = n;          // deduces int&
auto&& second = 5;         // deduces int&&
static_assert(std::is_same_v<decltype(first), int&>);
static_assert(std::is_same_v<decltype(second), int&&>);
\end{lstlisting}

W typie wyniku funkcji zapis \texttt{decltype(auto)} wymaga
ostrożności. Instrukcja \texttt{return (x);} może zwrócić
referencję do zmiennej lokalnej, która przestanie istnieć.
Zwykłe \texttt{auto} jest często bezpieczniejszym wyborem,
gdy funkcja ma zwrócić własny obiekt przez wartość.

\section{Referencje lvalue i rvalue}
\label{sec:move-references}

\begin{lstlisting}
#include <string>

std::string name = "Ada";
std::string& a = name;            // binds to mutable lvalue
const std::string& b = name;      // binds to lvalue
const std::string& c = "temporary"; // can bind to temporary
std::string&& d = std::string{"x"}; // binds to rvalue

// std::string& e = std::string{"x"}; // error
// std::string&& f = name;           // error
\end{lstlisting}

Referencja \texttt{T\&} umożliwia pracę z istniejącym, zwykle
nazwanym obiektem; \texttt{T\&\&} może wiązać się z r-wartością.
\texttt{const T\&} wiąże się zarówno z lvalue, jak i z rvalue,
nie dając prawa do modyfikowania obiektu poprzez tę referencję.

\subsection{Wybór przeciążenia}

\begin{lstlisting}
#include <iostream>
#include <string>
#include <utility>

void inspect(const std::string&) {
    std::cout << "const lvalue reference\n";
}
void inspect(std::string&&) {
    std::cout << "rvalue reference\n";
}

int main() {
    std::string text = "hello";
    inspect(text);             // const lvalue reference
    inspect(std::string{"x"}); // rvalue reference
    inspect(std::move(text));  // rvalue reference
}
\end{lstlisting}

Istnienie przeciążenia \texttt{f(T\&\&)} samo w sobie nie przesądza
o rzeczywistym przeniesieniu czegokolwiek. Funkcja może tylko
odczytać argument; \emph{move} jest operacją wykonywaną przez
konstruktor lub przypisanie, a nie przez mechanizm wiązania referencji.

\subsection{Wydłużenie życia obiektu tymczasowego}

W typowych deklaracjach lokalnych związanie obiektu tymczasowego
bezpośrednio z referencją może wydłużyć jego życie do końca życia
tej referencji:
\begin{lstlisting}
#include <string>

const std::string& a = std::string{"safe until a ends"};
std::string&& b = std::string{"safe until b ends"};
\end{lstlisting}
Nie wolno jednak uogólniać tej reguły na zwracanie referencji ani
na dowolne przekazywanie referencji dalej. Przykładowo:
\begin{lstlisting}
#include <string>

const std::string& broken() {
    return std::string{"dangling"}; // dangling reference
}
\end{lstlisting}
Zwrócona referencja nie utrzymuje tego obiektu przy życiu.
Jeżeli funkcja ma \emph{wytworzyć} wynik, zwykle powinna zwracać go
przez wartość.

\section{Co naprawdę robi \texttt{std::move}?}
\label{sec:move-std-move}

Nazwa jest nieco myląca: \texttt{std::move(x)} \textbf{nie przenosi
zasobów samodzielnie}. W przybliżeniu wykonuje rzutowanie do
referencji rvalue odpowiedniego typu. Dopiero dalszy konstruktor,
operator przypisania lub inna funkcja może potraktować wynik
jako obiekt nadający się do przeniesienia.

\begin{lstlisting}
#include <type_traits>
#include <utility>

template<class T>
constexpr std::remove_reference_t<T>&&
my_move(T&& value) noexcept {
    return static_cast<std::remove_reference_t<T>&&>(value);
}
\end{lstlisting}

Powyższa implementacja ilustruje istotę działania
\texttt{std::move}. W rzeczywistym kodzie korzystamy z funkcji
standardowej z nagłówka \texttt{<utility>}.

\begin{lstlisting}
#include <string>
#include <utility>

std::string first = "long text";
std::string second = std::move(first);
// A move constructor is selected for second.
// first is valid; do not assume its exact contents.

first = "reused";             // valid: assignment
first.clear();                // valid: clear()
\end{lstlisting}

\paragraph{Stan obiektu po przeniesieniu.}
Typy biblioteczne zwykle gwarantują stan \emph{poprawny, lecz
nieokreślony co do zawartości} (\textit{valid but unspecified}), o ile
kontrakt konkretnego typu nie obiecuje więcej. Takiego obiektu
można bezpiecznie zniszczyć, przypisać mu nową wartość albo
wykonać operację, której wymagania wstępne na to pozwalają.
Nie należy jednak \emph{zakładać}, że każdy przeniesiony
\texttt{std::string} jest pusty. Dla niektórych typów, np.
\texttt{std::unique\_ptr}, stan po przeniesieniu określono dokładniej.

\subsection{Problem obiektu \texttt{const}}

\begin{lstlisting}
#include <string>
#include <utility>

const std::string original = "cannot be emptied here";
std::string other = std::move(original);
// Typically: copy, not move (const string&& cannot bind to
// string's non-const move constructor).
\end{lstlisting}

\texttt{std::move(original)} ma w tym przypadku typ
\texttt{const std::string\&\&}. Typowy konstruktor przenoszący
przyjmuje \texttt{std::string\&\&}, więc nie może otrzymać
argumentu \texttt{const}. Kompilator wybiera konstruktor kopiujący
przyjmujący \texttt{const std::string\&}. \textbf{Samo wystąpienie
\texttt{std::move} w kodzie nie gwarantuje przeniesienia}.\footnote{Dla
typów niekopiowalnych ten sam zapis może skończyć się błędem
kompilacji.}

\section{Konstruktor przenoszący i przenoszący operator przypisania}
\label{sec:move-special-members}

W klasie \texttt{T} konstruktor przenoszący ma zwykle sygnaturę
\texttt{T(T\&\&)}, natomiast przenoszący operator przypisania:
\texttt{T\& operator=(T\&\&)}. Nie należy mylić ich z kopiowaniem:
\begin{lstlisting}
struct Item {
    Item() = default;
    Item(const Item&);            // copy constructor
    Item& operator=(const Item&); // copy assignment
    Item(Item&&) noexcept;        // move constructor
    Item& operator=(Item&&) noexcept; // move assignment
    ~Item();                      // destructor
};
\end{lstlisting}

\subsection{Przykład z obserwacją operacji}
\label{sec:move-trace}

\begin{lstlisting}
#include <iostream>
#include <utility>

struct Trace {
    Trace() { std::cout << "default\n"; }
    Trace(const Trace&) { std::cout << "copy ctor\n"; }
    Trace(Trace&&) noexcept { std::cout << "move ctor\n"; }
    Trace& operator=(const Trace&) {
        std::cout << "copy assign\n";
        return *this;
    }
    Trace& operator=(Trace&&) noexcept {
        std::cout << "move assign\n";
        return *this;
    }
};

int main() {
    Trace a;                    // default
    Trace b = a;                // copy ctor
    Trace c = std::move(a);     // move ctor
    b = c;                      // copy assign
    b = std::move(c);           // move assign
}
\end{lstlisting}

Konstruktor tworzy \emph{nowy} obiekt, a operator przypisania
zmienia już istniejący. W tym przykładzie \texttt{Trace} jedynie
wypisuje komunikaty. W klasie zarządzającej zasobem konstruktor
przenoszący powinien przekazać własność zasobu, a operator
przypisania musi dodatkowo poprawnie pozbyć się zasobu
posiadanego wcześniej przez cel.

\subsection{Pełny przykład: klasa samodzielnie zarządzająca tablicą}
\label{sec:move-raw-buffer}

To przykład \emph{edukacyjny}. W rzeczywistym projekcie zwykle
lepiej użyć standardowego kontenera, np. \texttt{std::vector<int>}.

\begin{lstlisting}
#include <algorithm>
#include <cstddef>
#include <utility>

class RawBuffer {
    std::size_t size_ = 0;
    int* data_ = nullptr;

public:
    explicit RawBuffer(std::size_t n)
        : size_(n), data_(n ? new int[n]{} : nullptr) {}

    ~RawBuffer() { delete[] data_; }

    RawBuffer(const RawBuffer& other)
        : RawBuffer(other.size_) {
        if (size_ != 0)
            std::copy_n(other.data_, size_, data_);
    }

    RawBuffer& operator=(const RawBuffer& other) {
        if (this != &other) {
            RawBuffer copy(other); // may throw; *this unchanged
            swap(copy);
        }
        return *this;
    }

    RawBuffer(RawBuffer&& other) noexcept
        : size_(std::exchange(other.size_, 0)),
          data_(std::exchange(other.data_, nullptr)) {}

    RawBuffer& operator=(RawBuffer&& other) noexcept {
        if (this != &other) {
            delete[] data_;
            size_ = std::exchange(other.size_, 0);
            data_ = std::exchange(other.data_, nullptr);
        }
        return *this;
    }

    void swap(RawBuffer& other) noexcept {
        using std::swap;
        swap(size_, other.size_);
        swap(data_, other.data_);
    }

    std::size_t size() const noexcept { return size_; }
    int& operator[](std::size_t i) { return data_[i]; }
    const int& operator[](std::size_t i) const { return data_[i]; }
};
\end{lstlisting}

\texttt{std::exchange(obj, value)} zwraca starą wartość \texttt{obj}
i zastępuje ją wartością \texttt{value}. Konstruktor przenoszący
przejmuje więc wskaźnik i rozmiar, a źródło ustawia w jawnie
określony stan pusty. Operator przenoszącego przypisania
usuwa dotychczasową tablicę celu, przejmuje nową i zabezpiecza
przypadek \texttt{b = std::move(b)}. W implementacji kopiowania
widać z kolei \emph{silną gwarancję wyjątkową}: jeśli utworzenie
kopii nie powiedzie się, pierwotny obiekt pozostaje bez zmian.

\section{Reguła pięciu, reguła trzech i reguła zera}
\label{sec:move-rule-of-five}

Klasy zarządzające zasobem często muszą kontrolować pięć operacji:
(1) destruktor, (2) konstruktor kopiujący, (3) kopiujący operator
przypisania, (4) konstruktor przenoszący oraz (5) przenoszący
operator przypisania. To \emph{reguła pięciu}.

Historyczna \emph{reguła trzech} odnosiła się do destruktora,
konstruktora kopiującego i kopiującego operatora przypisania,
zanim język udostępnił move semantics.

\paragraph{Ważne: deklarowanie destruktora może blokować generowanie
konstruktora przenoszącego.}

\begin{lstlisting}
#include <memory>
#include <type_traits>

struct WithDestructor {
    std::unique_ptr<int> p;
    ~WithDestructor() = default;
};

static_assert(!std::is_move_constructible_v<WithDestructor>);

struct WithMove {
    std::unique_ptr<int> p;
    ~WithMove() = default;
    WithMove() = default;
    WithMove(WithMove&&) noexcept = default;
    WithMove& operator=(WithMove&&) noexcept = default;
};
static_assert(std::is_move_constructible_v<WithMove>);
\end{lstlisting}

W \texttt{WithDestructor} użytkownik zadeklarował destruktor,
więc kompilator nie deklaruje automatycznie operacji przenoszących.
Kopiowanie jest niedostępne z powodu \texttt{unique\_ptr}.
Samo \texttt{= default} przy destruktorze nie usuwa tej reguły.

\subsection{Najlepsza praktyka: reguła zera}

\emph{Reguła zera} mówi, że klasy, które nie są bezpośrednio
odpowiedzialne za posiadanie surowych zasobów, zwykle nie powinny
same definiować żadnej z pięciu operacji.

\begin{lstlisting}
#include <string>
#include <vector>

struct Dataset {
    std::string title;
    std::vector<double> values;
};

Dataset original{"results", {1.0, 2.0, 3.0}};
Dataset copy = original;           // deep copy of vector data
Dataset moved = std::move(original); // member-wise move
\end{lstlisting}

To kontenery standardowe odpowiadają za przydział i zwalnianie
pamięci. Automatycznie generowane operacje \texttt{Dataset}
wykonują odpowiednie operacje na jego składowych. \textbf{Zalecenie:}
zacznij od reguły zera; własny konstruktor przenoszący pisz
jedynie wtedy, gdy zachowanie domyślne nie jest właściwe.

\section{Przenoszenie a wyjątki: \texttt{noexcept}}
\label{sec:move-noexcept}

Kontenery takie jak \texttt{std::vector<T>} przy realokacji muszą
przenieść lub skopiować dotychczasowe elementy do nowego obszaru.
Jeśli przenoszenie \texttt{T} może rzucać wyjątki, a \texttt{T}
można kopiować, biblioteka może wybrać kopiowanie, by zachować
silniejsze gwarancje wyjątkowe. Właśnie dlatego poprawne operacje
przenoszące, które nie mogą rzucać wyjątków, warto deklarować
\texttt{noexcept}.

\begin{lstlisting}
#include <type_traits>
#include <utility>

struct Resource {
    Resource() = default;
    Resource(const Resource&) = default;
    Resource(Resource&&) noexcept = default;
};

static_assert(std::is_nothrow_move_constructible_v<Resource>);

// move_if_noexcept can provide an lvalue reference if
// moving might throw and a copy is available.
\end{lstlisting}

\texttt{std::move\_if\_noexcept(x)} zwraca odpowiednio
\texttt{T\&\&} albo \texttt{const T\&}, zależnie od własności
konstruktorów danego typu. To narzędzie używane przy implementowaniu
kontenerów, a nie uniwersalny zamiennik \texttt{std::move}.

\paragraph{Ostrożnie z \texttt{noexcept}.}
Wyjście wyjątku poza funkcję zadeklarowaną jako
\texttt{noexcept} wywołuje \texttt{std::terminate}.
Dlatego bezwyjątkowość musi wynikać z rzeczywistych operacji składowych.

\section{Zwijanie referencji (\textit{reference collapsing})}
\label{sec:move-ref-collapse}

Gdy referencje występują pośrednio poprzez aliasy lub
podstawienie argumentów szablonu, stosuje się \emph{zwijanie
referencji}. Wszystkie przypadki przedstawia tabela:

\begin{center}
\begin{tabular}{|c|c|c|}
\hline
\textbf{Typ pośredni} & \textbf{Zapis} & \textbf{Wynik} \\
\hline
\texttt{T\&} i \texttt{\&} & \texttt{T\& \&} & \texttt{T\&} \\
\texttt{T\&} i \texttt{\&\&} & \texttt{T\& \&\&} & \texttt{T\&} \\
\texttt{T\&\&} i \texttt{\&} & \texttt{T\&\& \&} & \texttt{T\&} \\
\texttt{T\&\&} i \texttt{\&\&} & \texttt{T\&\& \&\&} & \texttt{T\&\&} \\
\hline
\end{tabular}
\end{center}

Krótko: jeżeli \emph{choć jeden} składnik jest \texttt{\&},
wynik jest \texttt{\&}; tylko \texttt{\&\&} połączone
z \texttt{\&\&} daje \texttt{\&\&}.

\begin{lstlisting}
#include <type_traits>

using L = int&;
using R = int&&;

static_assert(std::is_same_v<L&, int&>);
static_assert(std::is_same_v<L&&, int&>);
static_assert(std::is_same_v<R&, int&>);
static_assert(std::is_same_v<R&&, int&&>);
\end{lstlisting}

Zwijanie referencji jest podstawą \emph{perfect forwarding}:
dzięki niemu jeden szablon może przyjmować zarówno lvalue,
jak i rvalue i przekazywać tę różnicę dalej.

\section{Referencje przekazujące i dedukcja \texttt{T\&\&}}
\label{sec:move-forwarding-ref}

Zapis \texttt{T\&\&} ma dwa znaczenia w zależności od kontekstu.
Jeśli \texttt{T} jest dedukowanym, niekwalifikowanym parametrem
typu \emph{tego samego} szablonu funkcji, \texttt{T\&\&}
to \emph{referencja przekazująca} (\textit{forwarding reference}).
W pozostałych typowych kontekstach jest to po prostu
referencja rvalue.

\begin{lstlisting}
#include <string>

template<class T>
void forwarding(T&& x);         // forwarding reference

template<class T>
void not_forwarding(const T&& x); // rvalue reference

void fixed(std::string&& x);    // rvalue reference

template<class T>
struct Holder {
    void receive(T&& x);        // NOT forwarding: T comes
                               // from the class template

    template<class U>
    void relay(U&& x);          // forwarding: U is deduced
};
\end{lstlisting}

Dla wywołania \texttt{forwarding(x)} z lvalue typu \texttt{int}
kompilator dedukuje \texttt{T = int\&}, zatem parametr ma po
zwinięciu typ \texttt{int\&}. Dla \texttt{forwarding(5)}
dedukuje \texttt{T = int}, a parametr ma typ \texttt{int\&\&}.

\begin{lstlisting}
#include <type_traits>

template<class T>
void check(T&& x) {
    if constexpr (std::is_lvalue_reference_v<T>) {
        // Called with an lvalue.
    } else {
        // Called with an rvalue.
    }
}
\end{lstlisting}

\paragraph{Drobne wyjątki.} \texttt{auto\&\&} zwykle również
pozwala na przekazywanie kategorii, lecz dedukcja \texttt{auto}
ma szczególne reguły dla list inicjalizacyjnych.
\texttt{const T\&\&} nie jest referencją przekazującą.

\section{Perfect forwarding i \texttt{std::forward}}
\label{sec:move-perfect-forwarding}

Przez \emph{perfect forwarding} rozumiemy przekazanie argumentu
do następnej funkcji przy zachowaniu kategorii wartości,
a także odpowiednich kwalifikatorów typu.

\begin{lstlisting}
#include <iostream>
#include <utility>

void target(int&)  { std::cout << "lvalue\n"; }
void target(int&&) { std::cout << "rvalue\n"; }

template<class T>
void incorrect(T&& value) {
    target(value); // value is a named variable: always lvalue
}

template<class T>
void correct(T&& value) {
    target(std::forward<T>(value));
}

int main() {
    int x = 1;
    incorrect(2); // lvalue (!)
    correct(x);   // lvalue
    correct(2);   // rvalue
}
\end{lstlisting}

W przybliżeniu \texttt{std::forward<T>(value)} zwraca
lvalue, jeżeli \texttt{T} wydedukowano jako referencję lvalue,
i xvalue w przeciwnym razie. \texttt{std::move(value)} zawsze
wymusza kategorię xvalue: w funkcji przekazującej byłoby to
błędem, ponieważ niepotrzebnie pozwalałoby na przeniesienie
zasobu z argumentu przekazanego jako lvalue.

\subsection{Przykład: fabryka obiektów}

\begin{lstlisting}
#include <memory>
#include <string>
#include <utility>

struct Person {
    std::string name;
    int age;

    Person(std::string n, int a)
        : name(std::move(n)), age(a) {}
};

template<class T, class... Args>
std::unique_ptr<T> make_object(Args&&... args) {
    return std::unique_ptr<T>(
        new T(std::forward<Args>(args)...)
    );
}

int main() {
    std::string first = "Ada";
    auto p = make_object<Person>(first, 37); // first copied
    auto q = make_object<Person>(std::move(first), 38);
}
\end{lstlisting}

To edukacyjna wersja fabryki. W produkcyjnym kodzie
\texttt{std::make\_unique<T>(...)} robi to samo w prostszy
sposób. W drugim wywołaniu użytkownik świadomie pozwolił na
przeniesienie \texttt{first}; w pierwszym pozostaje on
nienaruszony. \texttt{std::forward<Args>(args)...} rozwija pakiet
argumentów omówiony w rozdziale~4.

\subsection{Nie przekazuj tego samego zasobu dwa razy bez zastanowienia}

\begin{lstlisting}
#include <utility>

template<class T>
void bad_double_forward(T&& x) {
    consume(std::forward<T>(x));
    consume(std::forward<T>(x)); // may use moved-from object
}
\end{lstlisting}

Ten fragment jest \emph{schematem błędu}: zakładamy, że istnieje
funkcja \texttt{consume}. Jeśli pierwsze wywołanie wykorzysta
argument rvalue przez przeniesienie, drugie otrzyma już obiekt
po przeniesieniu. Referencję przekazującą traktuj zatem jako
przeznaczoną zwykle do \emph{jednego} przekazania dalej.

\section{Perfect forwarding a przeciążenia i szablony klas}

Konstruktor z referencją przekazującą potrafi być bardzo zachłanny
podczas rozstrzygania przeciążeń. Może np. wygrać z konstruktorem
kopiującym przy argumentach nie-\texttt{const}.

\begin{lstlisting}
#include <string>
#include <utility>

struct Named {
    std::string name;

    Named(const Named&) = default;

    template<class U>
    explicit Named(U&& u)
        : name(std::forward<U>(u)) {}
};

// Named a{"text"};
// Named b(a);  // may select the constructor template,
//              // then fail to initialize std::string from Named&
\end{lstlisting}

Zwykle lepiej zacząć od zwykłych przeciążeń albo konstruktora
przyjmującego konkretny typ. Jeżeli szablon jest potrzebny,
można ograniczyć go koncepcją:

\begin{lstlisting}
#include <concepts>
#include <string>
#include <type_traits>
#include <utility>

struct SafeNamed {
    std::string name;

    SafeNamed(const SafeNamed&) = default;
    SafeNamed(SafeNamed&&) = default;

    template<class U>
      requires (!std::same_as<std::remove_cvref_t<U>, SafeNamed>
                && std::constructible_from<std::string, U&&>)
    explicit SafeNamed(U&& u)
        : name(std::forward<U>(u)) {}
};
\end{lstlisting}

Takie powiązanie rozdziału o perfect forwarding z koncepcjami
z rozdziału~5 jest bardzo przydatne
w bibliotekach generycznych.

\section{Copy elision: kiedy nie ma ani kopii, ani przeniesienia?}
\label{sec:move-elision}

Prosty model \emph{zwrócenie przez wartość = kopia lub move}
jest niepoprawny w nowoczesnym C++. Od C++17 pewne wyrażenia
\texttt{prvalue} inicjalizują wynik \emph{bezpośrednio w miejscu
docelowym}. Nie istnieje dodatkowy obiekt, który trzeba by
skopiować lub przenieść.

\begin{lstlisting}
struct NonMovable {
    NonMovable() = default;
    NonMovable(const NonMovable&) = delete;
    NonMovable(NonMovable&&) = delete;
};

NonMovable create() {
    return NonMovable{}; // guaranteed elision, C++17+
}

int main() {
    NonMovable x = create(); // OK in C++17+
}
\end{lstlisting}

Nie jest to zwykła optymalizacja, którą kompilator \emph{może}
wyłączyć: w pokazanym przypadku wynika to z semantyki języka.

\subsection{NRVO a gwarantowana elizja}

\begin{lstlisting}
#include <string>

std::string first() {
    return std::string{"hello"}; // mandatory direct construction
}

std::string second() {
    std::string s = "hello";
    return s;                    // NRVO may occur
}
\end{lstlisting}

\textit{Named Return Value Optimization} (NRVO) dotyczy
zwracania nazwanej lokalnej zmiennej. W C++20 kompilator \emph{może}
pominąć utworzenie dodatkowego obiektu. Jeśli NRVO nie nastąpi,
wyrażenie \texttt{s} w instrukcji \texttt{return} może
uczestniczyć w niejawnej próbie przeniesienia.

\paragraph{Nie pisz automatycznie \texttt{return std::move(s);}.}

\begin{lstlisting}
std::string make_text() {
    std::string s = "result";
    return s;             // allows NRVO
}

std::string less_good() {
    std::string s = "result";
    return std::move(s);  // inhibits NRVO; move may be required
}
\end{lstlisting}

\texttt{std::move(s)} zmienia wynik w xvalue i tym samym
nie spełnia warunków NRVO dotyczących nazwy obiektu.
Jeśli funkcja zwraca własną lokalną zmienną przez wartość,
zazwyczaj najlepsze jest \texttt{return s;}.

\section{Argumenty funkcji: kiedy kopiować, kiedy przenosić?}
\label{sec:move-api}

Wybór sposobu przyjmowania argumentu jest częścią kontraktu API.
Nie chodzi jedynie o wydajność, ale przede wszystkim o własność
i intencję.

\begin{center}
\begin{tabular}{|p{3.4cm}|p{4.3cm}|p{4.7cm}|}
\hline
\textbf{Parametr} & \textbf{Znaczenie} & \textbf{Kiedy używać} \\
\hline
\texttt{const T\&} & Odczyt bez przejmowania własności. & Duże obiekty, gdy wystarczy tylko odczyt. \\
\hline
\texttt{T\&} & Modyfikacja obiektu należącego do wywołującego. & Parametry wyjściowe i operacje in-place. \\
\hline
\texttt{T} & Własna wartość; może być skopiowana albo przeniesiona. & Małe typy i funkcje mające zachować wartość. \\
\hline
\texttt{T\&\&} & Dostęp do argumentu będącego r-wartością. & Przeciążenia zależne od kategorii oraz operacje konsumujące. \\
\hline
\texttt{T\&\&} z dedukcją \texttt{T} & Referencja przekazująca. & Opakowania i fabryki zachowujące kategorię argumentów. \\
\hline
\end{tabular}
\end{center}

\subsection{Idiom \emph{pass by value, then move}}

\begin{lstlisting}
#include <string>
#include <utility>

class User {
    std::string name_;
public:
    explicit User(std::string name)
        : name_(std::move(name)) {}
};

std::string a = "Alice";
User first(a);                    // copy into parameter, then move
User second(std::string{"Bob"});  // parameter from prvalue, then move
\end{lstlisting}

To dobre rozwiązanie, gdy funkcja \emph{zawsze} potrzebuje własnej
kopii/przeniesienia argumentu. Nie jest jednak magicznie
najszybsze w każdym przypadku: dla lvalue może wymagać kopiowania
i jeszcze jednego przeniesienia, a duże obiekty niemożliwe do
przeniesienia mogą wymagać osobnego projektu API.

\subsection{Przekazywanie własności \texttt{unique\_ptr}}

\begin{lstlisting}
#include <memory>
#include <utility>

void adopt(std::unique_ptr<int> p) {
    // Ownership transferred into this function.
}

int main() {
    auto p = std::make_unique<int>(42);
    adopt(std::move(p));
    // p == nullptr after moving this unique_ptr
}
\end{lstlisting}

Przekazanie \texttt{unique\_ptr} przez wartość oznacza przejęcie
własności. Jeżeli funkcja ma jedynie korzystać z obiektu,
zwykle lepiej przyjąć \texttt{T\&} lub \texttt{const T\&}.
Temat własności zasobów omówiono w rozdziale~3.

\section{Rvalue-qualified methods i świadome udostępnianie zasobów}
\label{sec:move-ref-qualified}

C++ pozwala ograniczać metodę do lvalue albo rvalue odbiorcy,
dopisując \texttt{\&} lub \texttt{\&\&} po liście parametrów.
Można to wykorzystać np. do udostępnienia zawartości obiektu
bez kopiowania, gdy całe opakowanie jest tymczasowe.

\begin{lstlisting}
#include <string>
#include <utility>

class Message {
    std::string data_;
public:
    explicit Message(std::string s) : data_(std::move(s)) {}

    const std::string& text() const & noexcept {
        return data_;              // lvalue owner: borrow
    }
    std::string text() && noexcept {
        return std::move(data_);   // rvalue owner: consume
    }
};

Message m{"hello"};
const std::string& borrowed = m.text();
std::string owned = Message{"world"}.text();
\end{lstlisting}

Wariant \texttt{\&\&} może bezpiecznie oddać wartość przez
przeniesienie. Gdybyśmy zwracali w nim referencję do składowej
tymczasowego obiektu, powstałaby referencja wisząca.
Projektując takie API, warto rozważyć również wywołania na
\texttt{const} r-wartościach i ich wymagane zachowanie.

\section{Kontenery, iteratory przenoszące i wzorce użytkowe}

\subsection{Dodawanie elementów do \texttt{std::vector}}

\begin{lstlisting}
#include <string>
#include <utility>
#include <vector>

std::vector<std::string> words;
std::string word = "one";

words.push_back(word);             // copy from lvalue
words.push_back(std::move(word));  // move from lvalue cast to xvalue
words.emplace_back(3, 'x');       // construct "xxx" in place
\end{lstlisting}

Metoda \texttt{emplace\_back} konstruuje nowy element z podanych
argumentów. Nie musi być szybsza od \texttt{push\_back}.
Gotowy obiekt można przekazać z pomocą \texttt{std::move}.
Podczas realokacji kontener nadal może kopiować lub przenosić
wcześniej dodane elementy.

\subsection{Przenoszenie całego zakresu}

\begin{lstlisting}
#include <algorithm>
#include <iterator>
#include <string>
#include <vector>

std::vector<std::string> source{"a", "b"};
std::vector<std::string> dest;

std::move(source.begin(), source.end(),
          std::back_inserter(dest)); // algorithm std::move

// source still contains two valid strings, but their
// contents are not assumed.
\end{lstlisting}

Nazwa \texttt{std::move} występuje również dla algorytmu z nagłówka
\texttt{<algorithm>}. W przeciwieństwie do pojedynczego
\texttt{std::move(x)} z nagłówka \texttt{<utility>}, algorytm
przypisuje kolejno elementy z zakresu docelowego, korzystając
z możliwości ich przenoszenia.

\section{Typowe błędy i pułapki}
\label{sec:move-errors}

\subsection{Błąd 1: przekonanie, że \texttt{std::move} niszczy obiekt}

\texttt{std::move} jest rzutowaniem, a nie destruktorem. Obiekt
źródłowy nadal istnieje aż do końca własnego czasu życia.
Może zmienić stan, jeśli na jego podstawie wykonano
rzeczywistą operację przeniesienia.

\subsection{Błąd 2: \texttt{std::move} z argumentu przekazanego jako \texttt{const T\&}}

\begin{lstlisting}
#include <string>
#include <utility>

std::string ineffective(const std::string& x) {
    return std::move(x); // copies for std::string
}
\end{lstlisting}

Jeżeli funkcja potrzebuje zwrócić niezależną kopię, czytelniej
napisać \texttt{return x;}. Jeśli ma konsumować argument,
powinna otrzymać wartość lub odpowiednią modyfikowalną
referencję. Nie wolno usuwać \texttt{const} przez rzutowanie,
by wymusić przeniesienie z oryginalnie stałego obiektu.

\subsection{Błąd 3: użycie nazwanego \texttt{T\&\&} bez \texttt{std::move}}

\begin{lstlisting}
#include <memory>
#include <utility>

void take(std::unique_ptr<int>&& p) {
    // auto q = p;            // error: p is an lvalue
    auto q = std::move(p);    // explicit ownership transfer
}
\end{lstlisting}

\subsection{Błąd 4: przechowywanie referencji do argumentu tymczasowego}

\begin{lstlisting}
#include <string>
#include <string_view>

std::string_view dangerous() {
    return std::string{"temporary"}; // dangling view
}
\end{lstlisting}

\texttt{std::string\_view} jest widokiem, a nie właścicielem
pamięci. Semantyka przenoszenia nie rozwiązuje problemów
związanych z czasem życia. Gdy wynik funkcji ma być własnością
wywołującego, zwracaj \texttt{std::string} przez wartość.

\subsection{Błąd 5: przenoszenie nie gwarantuje zachowania wskaźników}

Przeniesienie kontenera może spowodować unieważnienie wcześniej
zachowanych iteratorów, referencji lub wskaźników zależnie od
typu kontenera, operacji i alokatora. Nie zakładaj ogólnej
zasady, że po \texttt{std::move} adresy danych zawsze pozostają
takie same. Sprawdź dokumentację odpowiedniego typu.

\subsection{Błąd 6: nadmierne stosowanie \texttt{std::forward}}

\texttt{std::forward<T>} ma sens głównie tam, gdzie \texttt{T}
pochodzi z dedukcji dla referencji przekazującej. Przy zwykłej
lokalnej zmiennej przeznaczonej do przekazania własności
należy używać \texttt{std::move}; przy zwracaniu lokalnej
zmiennej przez wartość zazwyczaj żadnego z nich.

\section{Przykład przekrojowy: kolejka zadań}
\label{sec:move-project}

Rozważmy prosty bufor zadań przechowujący funkcje do
późniejszego wykonania. Użyjemy \texttt{std::unique\_ptr},
aby zadanie jednoznacznie posiadało zasób. Poniższy przykład
jest niezależny od wielowątkowości; \emph{kolejka zadań}
nie uruchamia niczego w osobnym wątku.

\begin{lstlisting}
#include <functional>
#include <iostream>
#include <memory>
#include <utility>
#include <vector>

class Tasks {
    std::vector<std::function<void()>> items_;
public:
    template<class F>
    void add(F&& f) {
        items_.emplace_back(std::forward<F>(f));
    }

    void run() {
        for (auto& f : items_) f();
        items_.clear();
    }
};

int main() {
    Tasks tasks;
    auto value = std::make_shared<int>(10);
    tasks.add([value] { std::cout << *value << '\n'; });
    tasks.run();
}
\end{lstlisting}

\texttt{std::function<void()>} w C++20 wymaga, aby przechowywany
obiekt wywoływalny był kopiowalny. Dlatego w przykładzie lambda
przechowuje \texttt{shared\_ptr}. Gdybyśmy zamiast niego
przechwycili przez przeniesienie \texttt{unique\_ptr}, lambda
stałaby się niekopiowalna i taki obiekt nie spełniłby
wymagań \texttt{std::function} w C++20. C++23 wprowadza
\texttt{std::move\_only\_function}, która rozwiązuje ten problem.
To pokazuje, że prawidłowy wybór sposobu przenoszenia zależy
od kontraktów pozostałych typów.

\section{Praktyczna ściąga i zalecenia projektowe}

\begin{center}
\begin{tabular}{|p{4.0cm}|p{8.4cm}|}
\hline
\textbf{Sytuacja} & \textbf{Zalecenie} \\
\hline
Zwracam utworzony obiekt & Zwróć go przez wartość; korzystaj z elizji lub NRVO. \\
\hline
Przenoszę zasób z lokalnego obiektu & Użyj \texttt{std::move(x)} w miejscu przekazania własności. \\
\hline
Piszę wrapper o argumentach \texttt{Args\&\&...} & Wykorzystaj \texttt{std::forward<Args>(args)...}. \\
\hline
Klasa ma \texttt{vector}/\texttt{string} jako składowe & Preferuj regułę zera i generowane operacje. \\
\hline
Klasa posiada goły wskaźnik do alokacji & Preferuj zmianę projektu na RAII; jeśli zarządzasz nim ręcznie, rozważ wszystkie operacje specjalne. \\
\hline
Kontener przenosi elementy przy realokacji & Właściwy, bezpieczny \texttt{noexcept} w konstruktorze move pomaga w gwarancjach wyjątkowych. \\
\hline
Argument będzie tylko odczytywany & Zwykle \texttt{const T\&} dla dużego \texttt{T}, a małe typy przez wartość. \\
\hline
Funkcja przejmuje \texttt{unique\_ptr} & Przyjmij go przez wartość; wywołujący świadomie używa \texttt{std::move}. \\
\hline
\end{tabular}
\end{center}

\paragraph{Najważniejsze rozróżnienie.}
\begin{enumerate}
\item \texttt{std::move} wyraża zgodę na potraktowanie obiektu jako źródła przeniesienia, ale samo zasobów nie przenosi.
\item \texttt{std::forward<T>} zachowuje kategorię argumentu dedukowanego w szablonie.
\item Konstruktor przenoszący to zwykle operacja \texttt{T(T\&\&)}, ale może wykonywać dowolną operację zgodną z kontraktem typu.
\item \texttt{T\&\&} jako typ parametru nie zawsze jest referencją przekazującą.
\item Nazwany parametr typu \texttt{T\&\&} jest wyrażeniem \texttt{lvalue}.
\item Zwracanie przez wartość jest często najlepszym sposobem przekazywania wyniku.
\item Kiedy tylko się da, stosuj regułę zera zamiast własnych pięciu funkcji specjalnych.
\end{enumerate}

\paragraph{Pytania kontrolne.}
\begin{enumerate}
\item Dlaczego wywołanie \texttt{g(x)} w funkcji \texttt{f(T\&\& x)} traktuje \texttt{x} jako lvalue?
\item Kiedy \texttt{std::move(x)} kończy się kopiowaniem zamiast przeniesienia?
\item Czym różnią się \texttt{return T\{\};} i \texttt{return local;} w C++20?
\item Dlaczego przy powiększaniu \texttt{vector} konstruktor move z \texttt{noexcept} ma znaczenie?
\item Kiedy użyć \texttt{std::forward<T>(x)}, a kiedy \texttt{std::move(x)}?
\item Dlaczego metoda zwracająca \texttt{const std::string\&} nie powinna ujawniać referencji do składowej obiektu tymczasowego?
\end{enumerate}

\paragraph{Dalsza lektura.} \emph{cppreference}: \textit{Value categories},
\textit{Move constructors}, \textit{Reference declaration},
\textit{Copy elision}, \textit{std::move}, \textit{std::forward}.
